\section{AMC 10 10B 2022}
        \begin{enumerate}
        	\item (\textit{AMC 10 10B 2022}) Define $x\diamond y$ to be $|x-y|$ for all real numbers $x$ and $y.$ What is the value of 
 \[(1\diamond(2\diamond3))-((1\diamond2)\diamond3)?\]
$\textbf{(A)}\ {-}2 \qquad \textbf{(B)}\ {-}1 \qquad \textbf{(C)}\ 0 \qquad \textbf{(D)}\ 1 \qquad \textbf{(E)}\ 2$

 \textbf{Answer: }A



\textbf{Solution 1}

We have \begin{align}(1\diamond(2\diamond3))-((1\diamond2)\diamond3) &= |1-|2-3|| - ||1-2|-3| \\&= |1-1| - |1-3| \\&= 0-2 \\&= \boxed{\textbf{(A)}\ {-}2}.\end{align}~MRENTHUSIASM

 



\textbf{Solution 2}

Observe that the $\diamond$ function is simply the positive difference between two numbers. Thus, we evaluate: the difference between $2$ and $3$ is $1;$ the difference between $1$ and $1$ is $0;$ the difference between $1$ and $2$ is $1;$ the difference between $1$ and $3$ is $2;$ and finally, $0-2=\boxed{\textbf{(A)}\ {-}2}.$



	\item (\textit{AMC 10 10B 2022}) In rhombus $ABCD$, point $P$ lies on segment $\overline{AD}$ so that $\overline{BP}$ $\perp$ $\overline{AD}$, $AP = 3$, and $PD = 2$. What is the area of $ABCD$? (Note: The figure is not drawn to scale.)


 
                    \begin{figure}[H]
                        \centering
                        \includegraphics[width=0.5\textwidth]{Problem/AMC_10_10B_2022_2_Problem_Figure_1.png}
                    \end{figure}
                    
$\textbf{(A) }3\sqrt 5 \qquad \textbf{(B) }10 \qquad \textbf{(C) }6\sqrt 5 \qquad \textbf{(D) }20\qquad \textbf{(E) }25$

 \textbf{Answer: }D



\textbf{Solution}

Notice that $BC = AD = AP + PD = 3 + 2 = 5$.Now, because the question tells us $ABCD$ is a rhombus, $BC = AB = 5$. Thus, by Pythagorean Theorem on $APB$, $BP = 4$.Therefore, 
 \[[ABCD] = AD \times BP = 5 \times 4 = \boxed{\text{D 20}}\]~DUODUOLUO(First time using $\LaTeX$, may have some minor mistakes)

 


	\item (\textit{AMC 10 10B 2022}) How many three-digit positive integers have an odd number of even digits?


 $\textbf{(A) }150\qquad\textbf{(B) }250\qquad\textbf{(C) }350\qquad\textbf{(D) }450\qquad\textbf{(E) }550$

 \textbf{Answer: }D



\textbf{Solution 1}

We use simple case work to solve this problem.

 Case 1: even, even, even = $4 \cdot 5 \cdot 5 = 100$

 Case 2: even, odd, odd = $4 \cdot 5 \cdot 5 = 100$

 Case 3: odd, even, odd = $5 \cdot 5 \cdot 5 = 125$

 Case 4: odd, odd, even = $5 \cdot 5 \cdot 5 = 125$

 Simply sum up the cases to get your answer, $100 + 100 + 125 + 125 = \boxed{\textbf{(D)~}450}$.

 - Wesseywes7254

 



\textbf{Solution 2 (Bijection)}

We will show that the answer is $450$ by proving a bijection between the three digit integers that have an even number of even digits and the three digit integers that have an odd number of even digits. For every even number with an odd number of even digits, we increment the number's last digit by $1$, unless the last digit is $9$, in which case it becomes $0$. It is very easy to show that every number with an even number of even digits is mapped to every number with an odd number of even digits, and vice versa. Thus, the answer is half the number of three digit numbers, or $\boxed{\textbf{(D)~}450}$




\textbf{Solution 3}

There are 9 possibilities for the first choice of a three-digit number in the hundreds place, excluding $0$. Now we have a few cases:

 Case 1: The first digit is odd

 In this case, only $1$ even digit is permitted, since $3$ is impossible. The even digit can be either in the tens or the ones, so we have $5\cdot5\cdot2$ ways to place the even digit and odd digit.

 Case 2: The first digit is even

 In this case, either all digits can be even or the rest of the digits can be odd. That's $5\cdot5$ for both of them, yielding $5\cdot5\cdot2$ in total.

 Thus, both cases result in multiplying the $9$ by $5\cdot5\cdot2$, or $\boxed{\textbf{(D)~}450}$

 -ichang3765

 


	\item (\textit{AMC 10 10B 2022}) A donkey suffers an attack of hiccups and the first hiccup happens at $4:00$ one afternoon. Suppose that
the donkey hiccups regularly every $5$ seconds. At what time does the donkey’s $700$th hiccup occur?


 $\textbf{(A) }15 \text{ seconds after } 4:58$
$\textbf{(B) }20 \text{ seconds after } 4:58$
$\textbf{(C) }25 \text{ seconds after } 4:58$
$\textbf{(D) }30 \text{ seconds after } 4:58$
$\textbf{(E) }35 \text{ seconds after } 4:58$

 \textbf{Answer: }A



\textbf{Solution 1}

Since the donkey hiccupped the 1st hiccup at $4:00$, it hiccupped for $5 \cdot (700-1) = 3495$ seconds, which is $58$ minutes and $15$ seconds, so the answer is $\boxed{\textbf{(A) }15 \text{ seconds after } 4:58}$.

 



\textbf{Solution 2 (Faster)}

We see that the minute has already been determined.The donkey hiccups once every 5 seconds, or 12 times a minute. $700\equiv 4 \pmod{12}$, so the 700th hiccup happened on the same second as the 4th, which occurred on the $5(4-1)=15$th second. $\boxed{\textbf{(A) }15 \text{ seconds after } 4:58}$.

 



\textbf{Solution 3 (Another Fast Way)}

We can add $7\cdot500=3500$ and then minus $5$ seconds.
 \begin{align*} 4:00+3500\text{ seconds} &= 4:00+3600\text{ seconds }-100\text{ seconds} \\ &=4:00+1\text{ hour }-100\text{ seconds} \\ &=5:00-100\text{ seconds} \\ &=20 \text{ seconds after } 4:58. \end{align*}
Finally, we minus $5$ seconds giving $\boxed{\textbf{(A) }15 \text{ seconds after } 4:58}$.

 


	\item (\textit{AMC 10 10B 2022}) What is the value of 
 \[\frac{\left(1+\frac13\right)\left(1+\frac15\right)\left(1+\frac17\right)}{\sqrt{\left(1-\frac{1}{3^2}\right)\left(1-\frac{1}{5^2}\right)\left(1-\frac{1}{7^2}\right)}}?\]
$\textbf{(A)}\ \sqrt3 \qquad\textbf{(B)}\ 2 \qquad\textbf{(C)}\ \sqrt{15} \qquad\textbf{(D)}\ 4 \qquad\textbf{(E)}\ \sqrt{105}$

 \textbf{Answer: }B



\textbf{Solution 1 (Difference of Squares)}

We apply the difference of squares to the denominator, and then regroup factors:
 \begin{align*} \frac{\left(1+\frac13\right)\left(1+\frac15\right)\left(1+\frac17\right)}{\sqrt{\left(1-\frac{1}{3^2}\right)\left(1-\frac{1}{5^2}\right)\left(1-\frac{1}{7^2}\right)}} &= \frac{\left(1+\frac13\right)\left(1+\frac15\right)\left(1+\frac17\right)}{\sqrt{\left(1+\frac13\right)\left(1+\frac15\right)\left(1+\frac17\right)}\cdot\sqrt{\left(1-\frac13\right)\left(1-\frac15\right)\left(1-\frac17\right)}} \\ &= \frac{\sqrt{\left(1+\frac13\right)\left(1+\frac15\right)\left(1+\frac17\right)}}{\sqrt{\left(1-\frac13\right)\left(1-\frac15\right)\left(1-\frac17\right)}} \\ &= \frac{\sqrt{\frac43\cdot\frac65\cdot\frac87}}{\sqrt{\frac23\cdot\frac45\cdot\frac67}} \\ &= \frac{\sqrt{4\cdot6\cdot8}}{\sqrt{2\cdot4\cdot6}} \\ &= \frac{\sqrt8}{\sqrt2} \\ &= \boxed{\textbf{(B)}\ 2}. \end{align*}
~MRENTHUSIASM

 



\textbf{Solution 2 (Brute Force)}

Since these numbers are fairly small, we can use brute force as follows: 
 \[\frac{(1+\frac{1}{3})(1+\frac{1}{5})(1+\frac{1}{7})}{\sqrt{(1-\frac{1}{3^2})(1-\frac{1}{5^2})(1-\frac{1}{7^2})}} =\frac{\frac{4}{3}\cdot\frac{6}{5}\cdot\frac{8}{7}}{\sqrt{\frac{8}{9}\cdot\frac{24}{25}\cdot\frac{48}{49}}} =\frac{\frac{4\cdot6\cdot8}{3\cdot5\cdot7}}{\sqrt{\frac{(2^3)(2^3\cdot3^1)(2^4\cdot3^1)}{(3^2)(5^2)(7^2)}}} =\frac{\frac{64}{35}}{\frac{96}{105}}=\frac{64}{35}\cdot\frac{105}{96}=\boxed{\textbf{(B)}\ 2}.\]~not\_slay

 



\textbf{Solution 3 (Brute Force)}

This solution starts precisely like the one above. We simplify to get the following:

 
 \[\frac{(1+\frac{1}{3})(1+\frac{1}{5})(1+\frac{1}{7})}{\sqrt{(1-\frac{1}{3^2})(1-\frac{1}{5^2})(1-\frac{1}{7^2})}} = \frac{\frac{4\cdot6\cdot8}{3\cdot5\cdot7}}{\sqrt{\frac{(2^3)(2^3\cdot3^1)(2^4\cdot3^1)}{(3^2)(5^2)(7^2)}}}\]

 But now, we can get a nice simplification as shown: 
 \[\frac{\frac{4\cdot6\cdot8}{3\cdot5\cdot7}}{\sqrt{\frac{(2^3)(2^3\cdot3^1)(2^4\cdot3^1)}{(3^2)(5^2)(7^2)}}} = \dfrac{\frac{4\cdot6\cdot8}{3\cdot5\cdot7}}{\sqrt{\frac{2^{10} \cdot 3^{2}}{3^2\cdot 5^2\cdot 7^2}}} = \dfrac{\frac{4\cdot6\cdot8}{3\cdot5\cdot7}}{\frac{2^5 \cdot 3}{3\cdot5\cdot 7}}  =\dfrac{4\cdot6\cdot8}{3\cdot5\cdot7} \hspace{0.05 in} \cdot \hspace{0.05 in}\dfrac{3\cdot5\cdot 7}{2^5 \cdot 3}  =\dfrac{2^6\cdot 3}{2^5\cdot 3} = \boxed{\textbf{(B)}\ 2}.\]



	\item (\textit{AMC 10 10B 2022}) How many of the first ten numbers of the sequence $121, 11211, 1112111, \ldots$ are prime numbers?


 $\textbf{(A) } 0 \qquad \textbf{(B) }1 \qquad \textbf{(C) }2 \qquad \textbf{(D) }3 \qquad \textbf{(E) }4$

 \textbf{Answer: }A



\textbf{Solution 1 (Generalization)}

The $n$th term of this sequence is
 \[\sum_{k=n}^{2n}10^k + \sum_{k=0}^{n}10^k = 10^n\sum_{k=0}^{n}10^k + \sum_{k=0}^{n}10^k = \left(10^n+1\right)\sum_{k=0}^{n}10^k.\]It follows that the terms are
 \begin{align*} 121 &= 11\cdot11, \\ 11211 &= 101\cdot111, \\ 1112111 &= 1001\cdot1111, \\ & \ \vdots \end{align*}
Therefore, there are $\boxed{\textbf{(A) } 0}$ prime numbers in this sequence.

 



\textbf{Solution 2 (Detailed Explanation of Solution 1)}

Denote this sequence as $a_{n}$, then we can find that
 \begin{align*} a_{1} &= 121 = 10^2 + 2\cdot10 + 1 = (10^2 + 10) + (10 + 1), \\ a_{2} &= 11211 = (10^4 + 10^3 + 10^2) + (10^2 + 10 + 1), \\ a_{3} &= 1112111 = (10^6 + 10^5 + 10^4 + 10^3) + (10^3 + 10^2 + 10 + 1), \\ & \ \vdots \end{align*}
So, we can induct that the general term is
 \begin{align*} a_n &= (10^{2n} + 10^{2n-1} + \ldots + 10^{n+1} + 10^n) + (10^n + 10^{n-1} + \ldots +10 + 1) \\ &= 10^n\cdot(10^n + 10^{n-1} + \ldots +10 + 1) + (10^n + 10^{n-1} + \ldots +10 + 1) \\ &= \left(10^n+1\right)\sum_{k=0}^{n}10^k. \end{align*}
Therefore, there are $\boxed{\textbf{(A) } 0}$ prime numbers in this sequence.

 



\textbf{Solution 3 (Simple Sums)}

Observe how 
 \begin{align*} 121 &= 110 + 11, \\ 11211 &= 11100 + 111, \\ 1112111 &= 1111000 + 1111, \\ & \ \vdots \end{align*}
all take the form of 
 \[\underbrace{111\ldots}_{n+1}\underbrace{00\ldots}_{n} + \underbrace{111\ldots}_{n+1} = \underbrace{111\ldots}_{n+1}(10^{n} + 1).\]Factoring each of the sums, we have 
 \[11(10+1), 111(100+1), 1111(1000+1), \ldots\] respectively. With each number factored, there are $\boxed{\textbf{(A) } 0}$ primes in the set.

 



\textbf{Solution 4 (Educated Guess)}

Note that $121$ is divisible by $11$ and $11211$ is divisible by $3$. Because this is Problem 6 of the AMC 10, we assume we do not need to check two-digit prime divisibility or use obscure theorems. Therefore, the answer is $\boxed{\textbf{(A) } 0}.$



	\item (\textit{AMC 10 10B 2022}) For how many values of the constant $k$ will the polynomial $x^{2}+kx+36$ have two distinct integer roots?


 $\textbf{(A)}\ 6 \qquad\textbf{(B)}\ 8 \qquad\textbf{(C)}\ 9 \qquad\textbf{(D)}\ 14 \qquad\textbf{(E)}\ 16$

 \textbf{Answer: }B



\textbf{Solution 1}

Let $p$ and $q$ be the roots of $x^{2}+kx+36.$ By \href{https://artofproblemsolving.com/wiki/index.php?title=Vieta\%27s\_Formulas}{Vieta's Formulas}, we have $p+q=-k$ and $pq=36.$

 It follows that $p$ and $q$ must be distinct factors of $36.$ The possibilities of $\{p,q\}$ are 
 \[\pm\{1,36\},\pm\{2,18\},\pm\{3,12\},\pm\{4,9\}.\]Each unordered pair gives a unique value of $k.$ Therefore, there are $\boxed{\textbf{(B) }8}$ values of $k,$ corresponding to $\mp37,\mp20,\mp15,\mp13,$ respectively.

 



\textbf{Solution 2}

Note that $k$ must be an integer. Using the \href{https://artofproblemsolving.com/wiki/index.php?title=Quadratic\_formula}{quadratic formula}, $x=\frac{-k \pm \sqrt{k^2-144}}{2}.$ Since $4$ divides $144$ evenly, $k$ and $k^2-144$ have the same parity, so $x$ is an integer if and only if $k^2-144$ is a perfect square.

 Let $k^2-144=n^2.$ Then, $(k+n)(k-n)=144.$ Since $k$ is an integer and $144$ is even, $k+n$ and $k-n$ must both be even. Assuming that $k$ is positive, we get $5$ possible values of $k+n$, namely $2, 4, 8, 6, 12$, which will give distinct positive values of $k$, but $k+n=12$ gives $k+n=k-n$ and $n=0$, giving $2$ identical integer roots. Therefore, there are $4$ distinct positive values of $k.$ Multiplying that by $2$ to take the negative values into account, we get $4\cdot2=\boxed{\textbf{(B) }8}$ values of $k$.

 



\textbf{Solution 3 (Pythagorean Triples)}

Proceed similar to Solution 2 and deduce that the discriminant of $x^{2}+kx+36$ must be a perfect square greater than $0$ to satisfy all given conditions. Seeing something like $k^2-144$ might remind us of a right triangle, where $k$ is the hypotenuse, and $12$ is a leg. There are four ways we could have this: a $9$-$12$-$15$ triangle, a $12$-$16$-$20$ triangle, a $5$-$12$-$13$ triangle, and a $12$-$35$-$37$ triangle.

 Multiply by $2$ to account for negative $k$ values (since $k$ is being squared), and our answer is $\boxed{\textbf{(B) }8}$.

 



\textbf{Solution 4}

Since $36 = 2^2\cdot3^2$, that means there are $(2+1)(2+1) = 9$ possible factors of $36$. Since $6 \cdot 6$ violates the distinct root condition, subtract $1$ from $9$ to get $8$. Each sum is counted twice, and we count of those twice for negatives. This cancels out, so we get $\boxed{\textbf{(B) }8}$.

 



\textbf{Solution 5 (Discriminant and Factor Rainbow)}

Since to have two solutions, the discriminant must be greater than $0$, $k^2 - 144 > 0$ so k must be greater than $12$ or less than $-12$(since it's squared). Next, we know that when factoring, the two numbers must multiply to $36$. So you can create a factor rainbow(remember from elementary school) which has the pairs of integers that multiply to $36( 1, 2, 3, 4, 6, 6, 9, 12, 18, 36 )$ (you just look at the position and the one on the opposite side, for example $2$ is the second value, and $18$ is the second last value, and they multiply to 36). See that there are $4$ pairs of values which add up to greater than $12$($6+6$ is equal to $12,$ not greater), and we look at the negative values of them too(since k is squared) we get $\boxed{\textbf{(B) }8}$ total values.

 


	\item (\textit{AMC 10 10B 2022}) Consider the following $100$ sets of $10$ elements each:

 \begin{align*} &\{1,2,3,\ldots,10\}, \\ &\{11,12,13,\ldots,20\},\\ &\{21,22,23,\ldots,30\},\\ &\vdots\\ &\{991,992,993,\ldots,1000\}. \end{align*}

How many of these sets contain exactly two multiples of $7$?


 $\textbf{(A)}\ 40\qquad\textbf{(B)}\ 42\qquad\textbf{(C)}\ 43\qquad\textbf{(D)}\ 49\qquad\textbf{(E)}\ 50$

 \textbf{Answer: }B



\textbf{Solution 1 (Competition)}

There are $\left\lfloor \frac{1000}{7} \right\rfloor = 142$ numbers divisible by 7. We split these into 100 sets containing 10 numbers each, giving us 1.42 multiples of 7 per set. After the first set, the numbers come evenly, and we multiply 100 by $1.42 - 1 = \boxed{\textbf{(B)}\ 42}.$

 (Genuinely genius solution ~ MW)

 



\textbf{Solution 2 (Casework)}

We apply casework to this problem. The only sets that contain two multiples of seven are those for which:

 
            \begin{enumerate} 
        \item The multiples of $7$ are $1\pmod{10}$ and $8\pmod{10}.$ That is, the first and eighth elements of such sets are multiples of $7.$The first element is $1+10k$ for some integer $0\leq k\leq99.$ It is a multiple of $7$ when $k=2,9,16,\ldots,93.$

        \item The multiples of $7$ are $2\pmod{10}$ and $9\pmod{10}.$ That is, the second and ninth elements of such sets are multiples of $7.$The second element is $2+10k$ for some integer $0\leq k\leq99.$ It is a multiple of $7$ when $k=4,11,18,\ldots,95.$

        \item The multiples of $7$ are $3\pmod{10}$ and $0\pmod{10}.$ That is, the third and tenth elements of such sets are multiples of $7.$The third element is $3+10k$ for some integer $0\leq k\leq99.$ It is a multiple of $7$ when $k=6,13,20,\ldots,97.$

            \end{enumerate}
            Each case has $\left\lfloor\frac{100}{7}\right\rfloor=14$ sets. Therefore, the answer is $14\cdot3=\boxed{\textbf{(B)}\ 42}.$




\textbf{Solution 3 (Find A Pattern)}

We find a pattern. 
 \begin{align*} &\{1,2,3,\ldots,10\}, \\ &\{11,12,13,\ldots,20\},\\ &\{21,22,23,\ldots,30\},\\ &\vdots\\ &\{991,992,993,\ldots,1000\}. \end{align*}
We can figure out that the first set has $1$ multiple of $7$. The second set also has $1$ multiple of $7$. The third set has $2$ multiples of $7$. The fourth set has $1$ multiple of $7$. The fifth set has $2$ multiples of $7$. The sixth set has $1$ multiple of $7$. The seventh set has $1$ multiple of $7$. The eighth set has $2$ multiples of $7$. Disregarding the first set and then calculating this pattern further, we can see (reasonably) that it repeats for each $1$ sets.We see that the pattern for the number of multiples per $7$ sets (again, disregarding the first set) goes: $1,2,1,2,1,1,2.$ So, for every $7$ sets after the first, there are three sets with $2$ multiples of $7$. We calculate $\left\lfloor\frac{100-1}{7}\right\rfloor$ and multiply that by $3$. (We also disregard the remainder of $1$ since it doesn't add any extra sets with $2$ multiples of $7$.). We get $14\cdot3= \boxed{\textbf{(B) }42}$.

 



\textbf{Solution 4 (Fastest)}

Each set contains exactly $1$ or $2$ multiples of $7$ (Because it will need at least $2$ * $7$ = $14$ numbers in a group for $3$ or more multiples of $7$).

 There are $\dfrac{1000}{10}=100$ total sets and $\left\lfloor\dfrac{1000}{7}\right\rfloor = 142$ multiples of $7$. Thus, there are $142-100=\boxed{\textbf{(B) }42}$ sets with $2$ multiples of $7$.

 



\textbf{Solution 5 (Simple Counting, Similar to Solution 1)}

Consecutive multiples of $7$ must differ by $7$. So, if a set $\{\ldots1,\ldots2,\ldots3,(\ldots),\ldots8,\ldots9,\ldots0\}$ contains two multiples of $7$, they must end with the digits $1$ and $8$, $2$ and $9$, or $3$ and $0$. This reduces the problem to counting the amount of multiples of $7$ less than $1000$ that end with $1$, $2$, and $3$.

 The first multiple of $7$ that ends with $1$ is $21$. The next multiple that ends with $1$ occurs $70$ later, since that is the smallest multiple of $7$ we can add to $21$ without affecting the last digit. The greatest number of $70$'s we can add to $21$ while keeping it less than $1000$ is $13$, because $21 + 70(13) = 931$. Therefore, the set of multiples of $7$ less than $1000$ ending with $1$ is $\{21 + 70(0), 21+70(1),\ldots,21+70(13)\}$, meaning there are $14$ of these particular multiples. We can use the same reasoning to count the multiples of $7$ that end with $2$ and $3$. 

 The first multiple of $7$ that ends with $2$ is $42$. The greatest number of $70$'s we can add to $42$ here is also $13$, since $42 + 70(13) = 952$. The set of multiples of $7$ less than $1000$ ending with $2$ is $\{42 + 70(0), 42+70(1),\ldots,42+70(13)\}$, giving $14$ multiples.

 The first multiple of $7$ that ends with $3$ is $63$. The greatest number of $70$'s we can add to $63$ here is yet again $13$, since $63 + 70(13) = 973$. The set of multiples of $7$ less than $1000$ ending with $3$ is $\{63 + 70(0), 63+70(1),\ldots,63+70(13)\}$, giving another $14$ multiples.

 In total, there are $14 + 14 + 14 = 42$ of these multiples, and so $\boxed{\textbf{(B) }42}$ sets with two multiples of $7$.

 



\textbf{Solution 6 (Very Fast System of Equations)}

Let $a$ be the number of sets with $1$ multiple of $7$ and $b$ be the number of sets with $2$ multiples of $7$. Note that it is impossible for a set to have more than two multiples of $7$. 

 Since there are a total of $100$ sets, $a+b=100$. Also, since there are $\lfloor \frac{1000}{7}\rfloor = 142$ multiples of $7$ between $1$ and $1000$, we must have $a+2b=142$. 

 Solving the system of equations

 
 \[a+b=100\]
 \[a+2b=142\]

 for $b$ gives us the answer of $\boxed{\textbf{(B) }42}$.

 



\textbf{Solution 7 (reasonable count)}

(Similar to Solution 3, but a little more intuitive and less numbers.) Note that this system loops every cycle of length $70$, or $7$ such sets. From $1$ to $70$, the multiples of $7$ are $7$, $14$, $21$, $28$, $35$, $42$, $49$, $56$, $63$, and $70$; note that $(21,28),(42,49),$ and $(63,70)$ are in the same sets. Thus, for every $7$ sets, we have $3$ sets with exactly two multiples of $7$. We have $100$ sets, which is $98+2=7\cdot14+2$; the first $98$ sets contain $\dfrac37(7\cdot14)=42$ desired sets. The last two sets comprise the integers from $981$ to $999$; the multiples of $7$ here are $987$ and $994$. Neither of the two last sets contains two multiples of $7$, so our answer is simply $\boxed{\textbf{(D) }42}$. 

 



\textbf{Solution 8 (Pigeonhole Principle, very fast)}

Consider any consecutive seven sets beginning on one more than a multiple of $70$. There are guaranteed to be $10$ multiples of $7$ in these seven sets. Note that no set may have three multiples of $7$. This means that some three of the ten sets have $2$ multiples of seven. There are $14$ such groups of seven sets, and then the sets $\{981, 982, \dots, 989\}$ and $\{991, 992, \dots, \}$. These last two sets are congruent to $\{1, 2, \dots, 9\}$ and $\{11, 12, \dots, 19\}$ modulo 7, thus making it easy to check that each only has one multiple of $7$. Therefore, there are $14 \cdot 3 = \boxed{\textbf{(D) }42}$ sets with $2$ multiples of $7$. 

 


	\item (\textit{AMC 10 10B 2022}) The sum

 \[\frac{1}{2!}+\frac{2}{3!}+\frac{3}{4!}+\cdots+\frac{2021}{2022!}\] can be expressed as $a-\frac{1}{b!}$, where $a$ and $b$ are positive integers. What is $a+b$?


 $\textbf{(A)}\ 2020 \qquad\textbf{(B)}\ 2021 \qquad\textbf{(C)}\ 2022 \qquad\textbf{(D)}\ 2023 \qquad\textbf{(E)}\ 2024$

 \textbf{Answer: }D



\textbf{Solution 1}

For all positive integers $n,$ we have 
 \[\frac{n}{(n+1)!}=\frac{(n+1)-1}{(n+1)!}=\frac{1}{n!}-\frac{1}{(n+1)!}.\]Note that the original sum is a telescoping series:
 \begin{align*} \frac{1}{2!}+\frac{2}{3!}+\frac{3}{4!}+\cdots+\frac{2021}{2022!}  &= \left(\frac{1}{1!}-\frac{1}{2!}\right)+\left(\frac{1}{2!}-\frac{1}{3!}\right)+\left(\frac{1}{3!}-\frac{1}{4!}\right)+\cdots+\left(\frac{1}{2021!}-\frac{1}{2022!}\right) \\ &= \frac{1}{1!}+\left(-\frac{1}{2!}+\frac{1}{2!}-\frac{1}{3!}+\frac{1}{3!}-\frac{1}{4!}+\cdots+\frac{1}{2021!}\right)-\frac{1}{2022!} \\ &= 1-\frac{1}{2022!}. \end{align*}
Therefore, the answer is $1+2022=\boxed{\textbf{(D)}\ 2023}.$




\textbf{Solution 2}

We add $\frac{1}{2022!}$ to the original expression 
 \[\left(\frac{1}{2!}+\frac{2}{3!}+\frac{3}{4!}+\cdots+\frac{2021}{2022!}\right)+\frac{1}{2022!}=\left(\frac{1}{2!}+\frac{2}{3!}+\frac{3}{4!}+\dots+\frac{2020}{2021!}\right)+\frac{1}{2021!}.\]This sum clearly telescopes, thus we end up with $\left(\frac{1}{2!}+\frac{2}{3!}\right)+\frac{1}{3!}=\frac{2}{2!}=1$



	\item (\textit{AMC 10 10B 2022}) Camila writes down five positive integers. The unique mode of these integers is $2$ greater than their median, and the median is $2$ greater than their arithmetic mean. What is the least possible value for the mode?


 $\textbf{(A)}\ 5 \qquad\textbf{(B)}\ 7 \qquad\textbf{(C)}\ 9 \qquad\textbf{(D)}\ 11 \qquad\textbf{(E)}\ 13$

 \textbf{Answer: }D



\textbf{Solution 1 (Variables)}

Let $M$ be the median. It follows that the two largest integers are both $M+2.$

 Let $a$ and $b$ be the two smallest integers such that $a<b.$ The sorted list is 
 \[a,b,M,M+2,M+2.\]Since the median is $2$ greater than their arithmetic mean, we have $\frac{a+b+M+(M+2)+(M+2)}{5}+2=M,$ or 
 \[a+b+14=2M.\]Note that $a+b$ must be even. We minimize this sum so that the arithmetic mean, the median, and the unique mode are minimized. Let $a=1$ and $b=3,$ from which $M=9$ and $M+2=\boxed{\textbf{(D)}\ 11}.$




\textbf{Solution 2 (Elimination)}

We can also easily test all the answer choices. (This strategy is generally good to use for multiple-choice questions if you don't have a concrete method to proceed with!)

 For answer choice $\textbf{(A)},$ the mode is $5,$ the median is $3,$ and the arithmetic mean is $1.$ However, we can quickly see this doesn't work, as there are five integers, and they can't have an arithmetic mean of $1$ while having a mode of $5.$ 

 Trying answer choice $\textbf{(B)},$ the mode is $7,$ the median is $5,$ and the arithmetic mean is $3.$ From the arithmetic mean, we know that all the numbers have to sum to $15.$ We know three of the numbers: $\underline{\hspace{3mm}},\underline{\hspace{3mm}},5,7,7.$ This exceeds the sum of $15.$

 Now we try answer choice $\textbf{(C)}.$ The mode is $9,$ the median is $7,$ and the arithmetic mean is $5.$ From the arithmetic mean, we know that the list sums to $25.$ Three of the numbers are $\underline{\hspace{3mm}},\underline{\hspace{3mm}},7,9,9,$ which is exactly $25.$ However, our list needs positive integers, so this won't work.

 Since we were really close on answer choice $\textbf{(C)},$ we can intuitively feel that the answer is probably going to be $\textbf{(D)}.$ We can confirm this by creating a list that satisfies the problem and choose $\textbf{(D)}: 1,3,9,11,11.$

 So, our answer is $\boxed{\textbf{(D)}\ 11}.$



	\item (\textit{AMC 10 10B 2022}) All the high schools in a large school district are involved in a fundraiser selling T-shirts. Which of the choices below is logically equivalent to the statement "No school bigger than Euclid HS sold more T-shirts than Euclid HS"?


 $\textbf{(A) }$ All schools smaller than Euclid HS sold fewer T-shirts than Euclid HS.


 $\textbf{(B) }$ No school that sold more T-shirts than Euclid HS is bigger than Euclid HS.


 $\textbf{(C) }$ All schools bigger than Euclid HS sold fewer T-shirts than Euclid HS.


 $\textbf{(D) }$ All schools that sold fewer T-shirts than Euclid HS are smaller than Euclid HS.


 $\textbf{(E) }$ All schools smaller than Euclid HS sold more T-shirts than Euclid HS.


 
 \textbf{Answer: }B



\textbf{Solution 1}

Let $B$ denote a school that is bigger than Euclid HS, and $M$ denote a school that sold more T-shirts than Euclid HS.

 It follows that $\neg B$ denotes a school that is \textbf{not bigger than} Euclid HS, and $\neg M$ denotes a school that \textbf{did not sell more} T-shirts than Euclid HS.

 Converting everything to conditional statements (if-then form), the given statement becomes 
 \[B\implies\neg M.\]Its contrapositive is $M\implies\neg B,$ which is $\boxed{\textbf{(B)}}.$

 Note that "not bigger than" does not mean "smaller than", and "not selling more" does not mean "selling fewer". There is an equality case. Therefore, none of the other answer choices is equivalent to $B\implies\neg M.$




\textbf{Solution 2 (Elimination)}

Suppose we have five schools: Euclid HS with $50$ students and $10$ T-shirts sold.

 
            \begin{itemize}
        \item School $A$ with $51$ students and $10$ T-shirts sold.
            \end{itemize}
            
            \begin{itemize}
        \item School $B$ with $49$ students and $10$ T-shirts sold.
            \end{itemize}
            
            \begin{itemize}
        \item School $C$ with $49$ students and $9$ T-shirts sold.
            \end{itemize}
            
            \begin{itemize}
        \item School $D$ with $51$ students and $9$ T-shirts sold (This configuration is legal.)
            \end{itemize}
            Then, school $B$ rules out $\textbf{(A)}$, school $A$ rules out $\textbf{(C)}$, school $D$ rules out $\textbf{(D)}$, and school $C$ rules out $\textbf{(E)}$, leaving us with $\boxed{\textbf{(B)}}$ as the correct answer.

 



\textbf{Solution 3 (Logic/Common Sense to Eliminate)}

First, you can eliminate A because the question doesn't say anything about smaller schools.Next, you can notice that B seems right at first glance, so you can keep it in contention.Next, you can eliminate C because it says they sold fewer, while the problem only states they didn't sell more than Euclid HS, so it can be fewer or equal to the number of shirts.Next, you can eliminate D because for the same reason as A.Finally, you can eliminate E for the same reason as A.Since everything except B is eliminated, $\boxed{\textbf{(B)}}$ is the correct answer.

 


	\item (\textit{AMC 10 10B 2022}) A pair of fair $6$-sided dice is rolled $n$ times. What is the least value of $n$ such that the probability that the sum of the numbers face up on a roll equals $7$ at least once is greater than $\frac{1}{2}$?


 $\textbf{(A) } 2 \qquad \textbf{(B) } 3 \qquad \textbf{(C) } 4 \qquad \textbf{(D) } 5 \qquad \textbf{(E) } 6$

 \textbf{Answer: }C



\textbf{Solution 1 (Complement)}

Rolling a pair of fair $6$-sided dice, the probability of getting a sum of $7$ is $\frac16:$ Regardless what the first die shows, the second die has exactly one outcome to make the sum $7.$ We consider the complement: The probability of not getting a sum of $7$ is $1-\frac16=\frac56.$ Rolling the pair of dice $n$ times, the probability of getting a sum of $7$ at least once is $1-\left(\frac56\right)^n.$

 Therefore, we have $1-\left(\frac56\right)^n>\frac12,$ or 
 \[\left(\frac56\right)^n<\frac12.\] Since $\left(\frac56\right)^4<\frac12<\left(\frac56\right)^3,$ the least integer $n$ satisfying the inequality is $\boxed{\textbf{(C) } 4}.$




\textbf{Solution 2 (Fakesolve)}

Let's try the answer choices. We can quickly find that when we roll $3$ dice, either the first and second sum to $7$, the first and third sum to $7$, or the second and third sum to $7$. There are $6$ ways for the first and second dice to sum to $7$,  $6$ ways for the first and third to sum to $7$, and $6$ ways for the second and third dice to sum to $7$. However, we overcounted (but not by much) so we can assume that the answer is $\boxed {\textbf{(C) }4}$.

 



\textbf{Solution 3 (Fakesolve)}

We can start by figuring out what the probability is for each die to add up to $7$ if there is only $1$ roll. We can quickly see that the probability is $\frac16$, as there are $6$ ways to make $7$ from $2$ numbers on a die, and there are a total of $36$ ways to add $2$ numbers on a die. And since each time we roll the dice, we are adding to the probability, we can conclude that the total probability for rolling a sum of $7$ in $n$ rolls would be $\frac16$$n$. The smallest number that satisfies this is $\boxed {\textbf{(C) }4}$.

 



\textbf{Solution 4 (Fakesolve)}

On each roll, there is a $\frac 16$ chance of rolling a sum of $7$. You would need $4$ of these rolls to get $4 \cdot \frac 16,$ which is larger than $\frac 12.$ Therefore, the answer is $\boxed {\textbf{(C) }4}.$



	\item (\textit{AMC 10 10B 2022}) The positive difference between a pair of primes is equal to $2$, and the positive difference between the cubes of the two primes is $31106$. What is the sum of the digits of the least prime that is greater than those two primes?


 $\textbf{(A)}\ 8 \qquad\textbf{(B)}\ 10 \qquad\textbf{(C)}\ 11 \qquad\textbf{(D)}\ 13 \qquad\textbf{(E)}\ 16$

 \textbf{Answer: }E



\textbf{Solution 1 (Simplest)}

Denote the two primes as \( a \) and \( b \). Then,

 \( a - b = 2 \)

 \( a^3 - b^3 = 31106 \)

 We see that \( a = 2 + b \)

 Now, we have \( (2+b)^3 - b^3 = 31106 \)

 We apply the binomial theorem (or just expand) to \( (2+b)^3 \), getting

 \( 8 + 12b + 6b^2 + b^3 - b^3 = 31106 \)

 \( \Rightarrow 8 + 12b + 6b^2 = 31106 \)

 Subtracting by 8 on both sides results in

 \( 6b^2 + 12b = 31098 \)

 Then dividing it all by 6 and pushing it all to one side results in

 \( b^2 + 2b - 5183 = 0\)

 \( (b-71)(b+73) = 0\)

 We see that \( b \in \{-73, 71\} \). We negate all negative values, and see that \( b = 71 \).

 Therefore \( a - b = 2 \), \( a = 2 + b \), \( a = 2 + 71 = 73 \).

 The next prime number greater than both of these is $79$, and therefore our answer is \( 7 + 9 = \) $\boxed{\textbf{(E) }16}$.

 



\textbf{Solution 2}

Let the two primes be $a$ and $b$ with $a>b$. We would have $a-b=2$ and $a^{3}-b^{3}=31106$. Using difference of cubes, we would have $(a-b)(a^{2}+ab+b^{2})=31106$. Since we know $a-b$ is equal to $2$, $(a-b)(a^{2}+ab+b^{2})$ would become $2(a^{2}+ab+b^{2})=31106$. Simplifying more, we would get $a^{2}+ab+b^{2}=15553$.

 Now let's introduce another variable. Instead of using $a$ and $b$, we can express the primes as $x+2$ and $x$ where $a$ is $x+2$ and b is $x$. Plugging $x$ and $x+2$ in, we would have $(x+2)^{2}+x(x+2)+x^{2}$. When we expand the parenthesis, it would become $x^{2}+4x+4+x^{2}+2x+x^{2}$. Then we combine like terms to get $3x^{2}+6x+4$ which equals $15553$. Then we subtract 4 from both sides to get $3x^{2}+6x=15549$. Since all three numbers are divisible by 3, we can divide by 3 to get $x^{2}+2x=5183$.

 Notice how if we add 1 to both sides, the left side would become a perfect square trinomial: $x^{2}+2x+1=5184$ which is $(x+1)^{2}=5184$. Since $2$ is too small to be a valid number, the two primes must be odd, therefore $x+1$ is the number in the middle of them. Conveniently enough, $5184=72^{2}$ so the two numbers are $71$ and $73$. The next prime number is $79$, and $7+9=16$ so the answer is  $\boxed{\textbf{(E) }16}$.

 A note: If you aren't entirely familiar with square trinomials, after we have $x^{2}+2x=5183$, we can simply guess and check from here. We know that $70^{2} = 4900$, and we notice that for $x^{2}+2x=5183$, $x$ must be a odd number. Thus we can guess that $x=71$, which proves to be right. Continuing, we then know the larger prime is 73, and we know that 75 and 77 aren't primes, so thus our next bigger prime is 79, and 7+9=16. Thus the answer is $\boxed{\textbf{(E) }16}$.

 Understanding the "Completing a Square" formula: Adding $1$ to $x^{2}+2x$ is called completing a square. In the expression $x^{2}+bx$, where $b$ is a constant, you can add $b/{2}$ to make the square $(x+b/{2})^{2}$.

 



\textbf{Solution 3}

Let the two primes be $a$ and $b$, with $a$ being the larger prime. We have $a - b = 2$, and $a^3 - b^3 = 31106$. Using difference of cubes, we obtain $a^2 + ab + b^2 = 15553$. Now, we use the equation $a - b = 2$ to obtain $a^2 - 2ab + b^2 = 4$. Hence, 
 \[a^2 + ab + b^2 - (a^2 - 2ab + b^2) = 3ab = 15553 - 4 = 15549\]
 \[ab = 5183.\]Because we have $a = b+2$, $ab = (b+1)^2 - (1)^2$. Thus, $(b+1)^2 = 5183 + 1 = 5184$, so $b+1 = 72$. This implies $a = 73$, $b = 71$, and thus the next biggest prime is $79$, so our answer is $7 + 9 = \boxed{\textbf{(E) }16}$




\textbf{Solution 4 (Estimation)}

Let the two primes be $p$ and $q$ such that $p-q=2$ and $p^{3}-q^{3}=31106$

 By the difference of cubes formula, $p^{3}-q^{3}=(p-q)(p^{2}+pq+q^{2})$

 Plugging in $p-q=2$ and $p^{3}-q^{3}=31106$,

 $31106=2(p^{2}+pq+q^{2})$

 Through the givens, we can see that $p \approx q$.

 Thus, $31106=2(p^{2}+pq+q^{2})\approx 6p^{2}\\p^2\approx \frac{31106}{6}\approx 5200$

 Recall that $70^2=4900$ and $80^2=6400$. It follows that our primes must be only marginally larger than $70$, where we conveniently find $p=73, q=71$

 The least prime greater than these two primes is $79 \implies 7 + 9 \implies \boxed{\textbf{(E) }16}$




\textbf{Solution 5}

Let the two primes be $x + 1$ and $x - 1$. Then, plugging it into the second condition, we get $(x + 1)^3 - (x - 1)^3 = 31106.$ Expanding the left side, 
 \[6x^2 + 2 = 31106 \implies x^2 = 5184.\] Taking the square root of both sides, we get that $x = 72$ and the larger prime is $73$. The smallest prime larger than $73$ is $79$, which has a digit sum of $7 + 9 = \boxed{16}.$

 - NL008

 



\textbf{Solution 6}

Let the two primes be $x$ and $y$. We write $x^3-y^3=31106$. So $(x-y)(x^2+xy+y^2)=31106$. 

 We manipulate the equation by subtracting $3xy$ to get a perfect square.

 $(x^2-2xy+y^2)=15553-3xy$
$(x-y)^2=15553-3xy$
$4=15553-3xy$
$3xy=15549$
$xy=5183$.

 We write out the factors of $5183$ and find $71$ and $73$.

 We can then list out primes greater than $73$ to find $79$, which has a digit sum of $\boxed{16}$.

 


	\item (\textit{AMC 10 10B 2022}) Suppose that $S$ is a subset of $\left\{ 1, 2, 3, \ldots , 25 \right\}$ such that the sum of any two (not necessarily distinct) elements of $S$ is never an element of $S.$ What is the maximum number of elements $S$ may contain?


 $\textbf{(A)}\ 12 \qquad\textbf{(B)}\ 13 \qquad\textbf{(C)}\ 14 \qquad\textbf{(D)}\ 15 \qquad\textbf{(E)}\ 16$

 \textbf{Answer: }B



\textbf{Solution 1 (Pigeonhole Principle)}

Let $M$ be the largest number in $S$.We categorize numbers $\left\{ 1, 2, \ldots , M-1 \right\}$ (except $\frac{M}{2}$ if $M$ is even) into $\left\lfloor \frac{M-1}{2} \right\rfloor$ groups, such that the $i$th group contains two numbers $i$ and $M-i$.

 Recall that $M \in S$ and the sum of two numbers in $S$ cannot be equal to $M$, and the sum of numbers in each group above is equal to $S$. Thus, each of the above $\left\lfloor \frac{M-1}{2} \right\rfloor$ groups can have at most one number in $S$.Therefore,
 \begin{align*} |S| & \leq 1 + \left\lfloor \frac{M-1}{2} \right\rfloor \\ & \leq 1 + \left\lfloor \frac{25}{2} \right\rfloor \\ & = 13. \end{align*}


 Next, we construct an instance of $S$ with $|S| = 13$.Let $S = \left\{ 13, 14, \ldots , 25 \right\}$.Thus, this set is feasible.Therefore, the most number of elements in $S$ is$\boxed{\textbf{(B) }13}$.

 



\textbf{Solution 2 (Pigeonhole v2)}

We construct a possible subset $S$ with $13$ elements by including all odd integers from $1$ to $25$, inclusive. $S=\left\{ 1, 3, 5, \cdots , 25 \right\}$. The sum of any $2$ elements is even, and thus cannot be an element of $S$.

 To show that $S$ cannot have more than $13$ elements, assume for sake of contradiction that $|S| \geq 14$. Let $S=\left\{ x_1, x_2, \cdots , x_n \right\}$ where $n \geq 14$ and $x_1 < x_2 < \cdots < x_n$. Because the sums of any $2$ (not necessarily distinct) elements do not appear in $S$, $x_1+x_i$ is not an element of $S$ for all $1 \leq i \leq n$. So, $x_1, x_2, \cdots , x_n , x_1+x_1, x_1+x_2, \cdots , x_1+x_n$ are all distinct integers. Let these integers be elements of the set $T$. $|T|=2n$, and because $n \geq 14$, $|T| \geq 28$. But all elements of $T$ must be $\geq x_1$ and $\leq x_1+x_n \leq x_1+25$, leaving only 26 possible values for the elements in $T$. By the Pigeonhole Principle, the elements cannot be distinct, and we have a contradiction.

 Thus, $\boxed{\textbf{(B) }13}$ is the maximum possible size of $S$.

 



\textbf{Solution 3 (Intuitive and Fast)}

Notice how \( 24 + 25 > 25 \), \( 23 + 24 > 25 \), and this continues backward until \( n + n+1 > 25 \). We then see that \( 2n+1 > 25 \rightarrow 2n > 24 \rightarrow n > 12 \).

 Thus the elements in the set must have the property \( n > 12 : n \in \mathbb{Z}^+ \) and therefore our set is \( \{13, 14, 15, \dots, 25\} \) in which there is \( 25 - 13 + 1 \Rightarrow \) $\boxed{\textbf{(B) }13}$ elements.

 


	\item (\textit{AMC 10 10B 2022}) Let $S_n$ be the sum of the first $n$ terms of an arithmetic sequence that has a common difference of $2$. The quotient $\frac{S_{3n}}{S_n}$ does not depend on $n$. What is $S_{20}$?


 $\textbf{(A) } 340 \qquad \textbf{(B) } 360 \qquad \textbf{(C) } 380 \qquad \textbf{(D) } 400 \qquad \textbf{(E) } 420$

 \textbf{Answer: }D



\textbf{Solution 1}

Let's say that our sequence is 
 \[a, a+2, a+4, a+6, a+8, a+10, \ldots.\]Then, since the value of n doesn't matter in the quotient $\frac{S_{3n}}{S_n}$, we can say that
 \[\frac{S_{3}}{S_1} = \frac{S_{6}}{S_2}.\]Simplifying, we get $\frac{3a+6}{a}=\frac{6a+30}{2a+2}$, from which 
 \[\frac{3a+6}{a}=\frac{3a+15}{a+1}.\] 
 \[3a^2+9a+6=3a^2+15a\] 
 \[6a=6\] Solving for $a$, we get that $a=1$. 

 Since the sum of the first $n$ odd numbers is $n^2$, $S_{20} = 20^2 = \boxed{\textbf{(D) } 400}$.

 Note: you could also plug in the formulas for $\frac{S_{3n}}{S_{n}}$ and simplify, getting $3+ \frac{6n}{a+n-1}$ You would then find $a = 1$~megacleverstarfish15

 
(Optimal competition solution if you're short on time.~Lunar Paradox)

 



\textbf{Solution 2 (Quick Insight)}

Recall that the sum of the first $n$ odd numbers is $n^2$.

 Since $\frac{S_{3n}}{S_{n}} = \frac{9n^2}{n^2} = 9$, we have $S_{20} = 20^2 = \boxed{\textbf{(D) } 400}$.

 



\textbf{Solution 3 (I didn't know Solution 1!)}

We have a slightly challenging problem :-(, but that's okay!

 $\textbf{If you didn't come up with the Solution 1 intuition}$, then here is a more direct approach.

 We want \(\frac{S_{3n}}{S_n}\) to be a natural (\(\frac{S_{3n}}{S_n} > 0\); basically a whole number \(\neq 0\)) number. Then only can the progression be incremental by 2.

 Assume that \(a_1 = 1\). Then, \(a_2 = 3\), \(a_3 = 5\), etc. We take the first term \(n=1\), which is 1. Then \(3n\); the next 3 terms will sum to 9. \(\frac{9}{1} = 9\), and this is an arithmetic sequence, and therefore works!

 What about \(a_1 = 0\). We immediately see that this would be undefined, so we cannot have \(a_1 = 0\). So we say \(a_1 = 2\). Then the sum for \(n=1\) is  2, and for \(3n\); it is simply 6. This is \(6/2 = 3\), so it works!

 Then, what about \(a_2 = 3\)? Then this means that the \(n=2\) sum is 4, and the \(3n\) sum is 36. This is \(36/4 = 9\), and this is the same as the difference before, proving that \(a_1 = 1\) is indeed correct.

 And then? What about \(a_2 = 4\)? This means \(n=2\) sum is 6, and the \(3n\) sum is 42, but uh oh, the progression breaks! Therefore, the evens cannot work.

 Therefore, \(a_1 = 1\), \(a_{20}\) is just the sum of the first 20 odd numbers, which is $\boxed{\textbf{(D) } 400}$

 \textbf{Remark 3.0.1} Hi! It me (Pinotation) again!

 If you are confused about why \( a_1 \neq 3 \) or \( a_1 \neq 4 \) or something along those lines, we can prove by induction.

 We have
 \[S_n=\frac{n}{2}\big(2a_1+(n-1)\cdot 2\big)=n(a_1+n-1).\]Then
 \[S_{3n}=3n(a_1+3n-1),\]so
 \[\frac{S_{3n}}{S_n}=\frac{3(a_1+3n-1)}{a_1+n-1}.\]

 Base case \( n=1 \):
 \[\frac{S_3}{S_1}=\frac{3(a_1+2)}{a_1}.\]

 Suppose for some \( n \) the ratio is independent of \( n \) and equal to a constant \( k \). Then
 \[\frac{S_{3n}}{S_n}=\frac{3(a_1+3n-1)}{a_1+n-1}=k.\]For \( n+1 \) we must also have
 \[\frac{S_{3(n+1)}}{S_{n+1}}=\frac{3(a_1+3n+2)}{a_1+n}=k.\]So
 \[\frac{3(a_1+3n-1)}{a_1+n-1}=\frac{3(a_1+3n+2)}{a_1+n}.\]Cross multiplying gives
 \[(a_1+3n-1)(a_1+n)=(a_1+3n+2)(a_1+n-1).\]Expanding,
 \[a_1^2+4na_1+a_1- n -3 = a_1^2+4na_1+a_1+2n-2.\]Simplifying,
 \[- n -3 = 2n -2,\]so
 \[3n = -1.\]This is impossible for integer \( n \). The only way the equality can hold for all \( n \) is if the proportionality factor between numerator and denominator is fixed. Comparing coefficients in
 \[3(a_1+3n-1)=C(a_1+n-1),\]we get \( C=9 \) and hence
 \[3a_1-3=9a_1-9 \implies a_1=1.\]

 I hope this helped! If not, then you can check Solution 1, as it is practically the friendlier version LOL!

 



\textbf{Solution 4 (Answer Choices)}

Let $a$ be the first element. Then, $S_{20}=20(a+19)$. Let it equal to all the answer choices respectively, we get $a=-2, -1, 0, 1, 2$. We test $\frac{S_{3n}}{S_n}$ for each $a$, using $n=1, 2$. We can see that for $a=-2, -1, 0$ there exists $n$ to let $S_n=0$, which implies they are invalid. For $a=2$ we can also see it's invalid, since $\frac{S_3}{S_1}=6\neq7=\frac{S_6}{S_2}$. Therefore, $a=1$, and $\boxed{\textbf{(D) } 400}$ is the correct choice.

 



\textbf{Solution 5 (No usage of initial term or difference)}

Let a and d be the initial term and difference respectively. (These won't be used it later)

 Notice that the terms $n+1,n+2,...,2n=(a+d+nd),(a+2d+nd),...,(a+nd+nd)=S_{n}+n(nd)=S_{n}+n^2d$. In a similar fashion, we have $S_{3n}=3S_{n}+2n^2+4n^2=3S_{n}+6n^2$.

 We are given that the quotient $\frac{S_{3n}}{S_n}$ does not depend on $n$. With the result above, we have  
 \[\frac{3S_{n}+6n^2}{S_n}\]where we get
 \[3+\frac{6n^2}{S_n}\]With the information above, we must have
 \[S_{n}=kn^2\]If $k\neq1$, then ${d\neq2}$, so we must have $S_{n}=n^2$.

 Hence, $S_{20}=20^2=\boxed{\textbf{(D) } 400}$



	\item (\textit{AMC 10 10B 2022}) The diagram below shows a rectangle with side lengths $4$ and $8$ and a square with side length $5$. Three vertices of the square lie on three different sides of the rectangle, as shown. What is the area of the region inside both the square and the rectangle?


 
                    \begin{figure}[H]
                        \centering
                        \includegraphics[width=0.5\textwidth]{Problem/AMC_10_10B_2022_16_Problem_Figure_1.png}
                    \end{figure}
                    
$\textbf{(A) }15\dfrac{1}{8}  \qquad \textbf{(B) }15\dfrac{3}{8}  \qquad \textbf{(C) }15\dfrac{1}{2}  \qquad \textbf{(D) }15\dfrac{5}{8}  \qquad \textbf{(E) }15\dfrac{7}{8}$

 \textbf{Answer: }D



\textbf{Solution 1}

Let us label the points on the diagram.

 
                    \begin{figure}[H]
                        \centering
                        \includegraphics[width=0.5\textwidth]{Solution/AMC_10_10B_2022_16_Solution_Figure_1.png}
                    \end{figure}
                    

 By doing some angle chasing using the fact that $\angle ACE$ and $\angle CEG$ are right angles, we find that $\angle BAC = \angle DCE = \angle FEG$. Similarly, $\angle ACB = \angle CED = \angle EGF$. Therefore, $\triangle ABC \sim \triangle CDE \sim \triangle EFG$.

 As we are given a rectangle and a square, $AB = 4$ and $AC = 5$. Therefore, $\triangle ABC$ is a $3$-$4$-$5$ right triangle and $BC = 3$. 

 $CE$ is also $5$. So, using the similar triangles, $CD = 4$ and $DE = 3$.

 $EF = DF - DE = 4 - 3 = 1$. Using the similar triangles again, $EF$ is $\frac14$ of the corresponding $AB$. So, 

 
 \begin{align*} [\triangle EFG] &= \left(\frac14\right)^2 \cdot [\triangle ABC] \\ &= \frac{1}{16} \cdot 6 \\ &= \frac38. \end{align*}


 Finally, we have 

 
 \begin{align*} [ACEG] &= [ABDF] - [\triangle ABC] - [\triangle CDE] - [\triangle EFG] \\ &= 7 \cdot 4 - \frac12 \cdot 3 \cdot 4 - \frac12 \cdot 3 \cdot 4 - \frac38 \\ &= 28 - 6 - 6 - \frac38 \\ &= \boxed{\textbf{(D) }15\dfrac{5}{8}}. \end{align*}





\textbf{Solution 2 (Clever)}

(Refer to the diagram above)Proceed the same way as Solution 1 until you get all of the side lengths. Then, it is clear that due to the answer choices, we only need to find the fractional part of the shaded area. The area of the whole rectangle is integral, as is the area of $\triangle ABC$, $\triangle CDE$, and the rectangle to the far left of the diagram. The area of $EFG$ is $\frac{3}{8}$ and thus the fractional part of the answer is $\frac{5}{8}$. The only answer choice that has $\frac{5}{8}$ in it is $\boxed{\textbf{(D) }15\dfrac{5}{8}}$




\textbf{Solution 3 (Intuitive Risking)}


                    \begin{figure}[H]
                        \centering
                        \includegraphics[width=0.5\textwidth]{Solution/AMC_10_10B_2022_16_Solution_Figure_2.png}
                    \end{figure}
                    We can quickly find out through Pythagorean Triples that \( GH = 3 \). Now, we make an educated guess that \( G \) is the midpoint of \( FI \), and therefore \( HI = 1 \).

 Noice! So, now what? Well, \( FG \) is then \( 4 \), and \( EF \) would then be \( 3 \).

 Given that info, we also see that \( DE = 1 \).

 Our area of the rectangle \( ADFI \) is \( 8 \times 4 = 32 \). We then see that we must subtract rectangle \( ABHI = 4 \times 1 = 4 \), triangle \( BHG = \frac{1}{2} \times 3 \times 4 = 6 \), and triangle \( GFE = \frac{1}{2} \times 4 \times 3 = 6 \). This gives \( 32 - (4 + 6 + 6) = 16 \).

 We make a guess that the intersection between \( B \) and \( D \), call it \( J \), to \( E \) is exactly \( 5/4 \). Then, through the Pythagorean theorem, \( JD = 1/2 \), and the area is \( \frac{1}{2} \times \frac{1}{2} \times 1 = \frac{1}{4} \). So, the area of that mini triangle is \( \frac{1}{4} \), in reality we know it is bigger. We subtract this from \( 16 \) to get \( 15 \frac{3}{4} \) and take the closest answer choice that is lower than the result we got, which is $\boxed{\textbf{(D) }15\dfrac{5}{8}}$.

 Take risks everyone! You might do great things like solve an AMC problem :D!

 



\textbf{Solution 4 (Coordinate Geometry)}

Same diagram as Solution 1, but added point $H$, which is $(4,7)$. I also renamed all the points to form coordinates using $B$ as the origin. 
                    \begin{figure}[H]
                        \centering
                        \includegraphics[width=0.5\textwidth]{Solution/AMC_10_10B_2022_16_Solution_Figure_3.png}
                    \end{figure}
                    

 In order to find the area, point $G$'s coordinates must be found. Notice how $EH$ and $AG$ intercept at point $G$. This means that we need to find the equations for $EH$ and $AG$ and make a system of linear equations. 

 Using the slope formula $m=\frac{y_{2} - y_{1}}{x_{2} - x_{1}}$, we get the slope for $EH$, which means $m=\frac{3-7}{7-4} = -\frac{4}{3}$

 Then, by using point slope form. $y-y_{1}=m(x-x_{1})$. We can say that the equation for $EH$ is $y-7=-\frac{4}{3}(x-4)$ or in this case, $y=-\frac{4}{3}x+12 \frac{1}{3}$.

 And it is easy to figure out that the equation for $AG$ is $y=4$.

 The best way to solve the system of linear equations is to substitute the $y$ for the $4$ in equation $EH$. $4=-\frac{4}{3}x+12 \frac{1}{3}$, so $x=6\frac{1}{4}$ and $y=4$ This would mean $G\left(6\frac{1}{4},4\right)$.

 Since we have our $G$ coordinate, we can continue with Solution 3, with the area of the trapezoid $\left(\frac{EG+AC}{2}\right)(CE)$, where $EG=\frac{5}{4}$ (using distance formula for $E$ to $G$), $AC=5$, and $CE=5$.

 By substitution, we get $\left(\frac{\frac{5}{4}+5}{2}\right)(5)=$$\boxed{\textbf{(D) }15\dfrac{5}{8}}$. 

 



\textbf{Solution 5 (FAST)}

Notice the small triangle in the upper right corner is a $3-4-5$ triangle. Then that triangle is similar to the big triangle by AA similarity. From that, do similar triangles and you find that the longer leg of the small triangle is 1. Then you find that the triangle below is $3-4-5$, so the side length of the rectangle (without the outer rectangle) is 7. afterwards you just add the half of the square + the remaining triangle which can be found by multiplying base and height (in which we already know)

 
I found the same solution when solving this, but essentially, we know all triangles in the picture are similar, and we know the lower left hand triangle is a $3-4-5$ triangle, and it is similar to the top triangle, where 4 corresponds to 5, therefore the unknown side length of that triangle is 15/4. Continuing on, the are of the white part of the square is $5 * 15/4 * 1/2 = 75/8,  25 - 75/8 = \fbox{(D)}$

 
-kaiser

 



\textbf{Solution 6 (Fastest Similar Triangles)}

For reference, use the points labelled in the diagram of Solution 3. Let the point one unit to the right of $A$ be $A'$ (so that $A'$ is one of the vertices of the square). The square means $A'C = 5$, so we get a $3$-$4$-$5$ triangle $A'BC$.

 $m\angle BA'C = 90 ^{\circ} - m\angle CA'G = m\angle GA'H$. 

 Therefore $\triangle GA'H$ is proportional to a $3$-$4$-$5$ triangle, with $HG$ corresponding to $3$ and $A'H$ corresponding to $4$. By similar triangles, we find 

 $HG = A'H \cdot \frac{3}{4} = \frac{15}{4}$.

 
            \begin{enumerate}[label=(\roman*)]
                <dt>then, finally, $[A'CEG] = [A'CEH] - [\triangle GA'H] = 5^2 - \frac{1}{2} \cdot \frac{15}{4} \cdot 5 = \boxed{\textbf{(D) }15\dfrac{5}{8}}.$</dt>
            \end{enumerate}
            


	\item (\textit{AMC 10 10B 2022}) One of the following numbers is not divisible by any prime number less than $10.$ Which is it?


 $\textbf{(A) } 2^{606}-1 \qquad\textbf{(B) } 2^{606}+1 \qquad\textbf{(C) } 2^{607}-1 \qquad\textbf{(D) } 2^{607}+1\qquad\textbf{(E) } 2^{607}+3^{607}$

 \textbf{Answer: }C



\textbf{Solution 1 (Modular Arithmetic)}

For $\textbf{(A)}$ modulo $3,$
 \begin{align*} 2^{606} - 1 & \equiv (-1)^{606} - 1 \\ & \equiv 1 - 1 \\ & \equiv 0 . \end{align*}
Thus, $2^{606} - 1$ is divisible by $3.$

 For $\textbf{(B)}$ modulo $5,$
 \begin{align*} 2^{606} + 1 & \equiv 2^{{\textrm Rem} ( 606, \phi(5) )} + 1 \\ & \equiv 2^{{\textrm Rem} ( 606, 4 )} + 1 \\ & \equiv 2^2 + 1 \\ & \equiv 0 . \end{align*}
Thus, $2^{606} + 1$ is divisible by $5.$

 For $\textbf{(D)}$ modulo $3,$
 \begin{align*} 2^{607} + 1 & \equiv (-1)^{607} + 1 \\ & \equiv - 1 + 1 \\ & \equiv 0 . \end{align*}
Thus, $2^{607} + 1$ is divisible by $3.$

 For $\textbf{(E)}$ modulo $5,$
 \begin{align*} 2^{607} + 3^{607} & \equiv 2^{607} + (-2)^{607} \\ & \equiv 2^{607} - 2^{607} \\ & \equiv 0 . \end{align*}
Thus, $2^{607} + 3^{607}$ is divisible by $5.$

 Therefore, the answer is $\boxed{\textbf{(C) }2^{607} - 1}.$




\textbf{Solution 2 (Factoring)}

We have
 \begin{alignat*}{8} 2\textasciicircum{}{606}-1 &= 4\textasciicircum{}{303}-1 &&= (4-1)(4\textasciicircum{}{302}+4\textasciicircum{}{301}+4\textasciicircum{}{300}+\cdots+4\textasciicircum{}0), \\  2\textasciicircum{}{606}+1 &= 4\textasciicircum{}{303}+1 &&= (4+1)(4\textasciicircum{}{302}-4\textasciicircum{}{301}+4\textasciicircum{}{300}-\cdots+4\textasciicircum{}0), \\  2\textasciicircum{}{607}+1 & &&= (2+1)(2\textasciicircum{}{606}-2\textasciicircum{}{605}+2\textasciicircum{}{604}-\cdots+2\textasciicircum{}0), \\ 2\textasciicircum{}{607}+3\textasciicircum{}{607} & &&= (2+3)(2\textasciicircum{}{606}\cdot3\textasciicircum{}0-2\textasciicircum{}{605}\cdot3\textasciicircum{}1+2\textasciicircum{}{604}\cdot3\textasciicircum{}2-\cdots+2\textasciicircum{}0\cdot3\textasciicircum{}{606}). \end{alignat*}We conclude that $\textbf{(A)}$ is divisible by $3$, $\textbf{(B)}$ is divisible by $5$,$\textbf{(D)}$ is divisible by $3$, and $\textbf{(E)}$ is divisible by $5$.

 Since all of the other choices have been eliminated, we are left with $\boxed{\textbf{(C) } 2^{607}-1}$.

 



\textbf{Solution 3 (Elimination)}

Mersenne Primes are primes of the form $2^n-1$, where $n$ is prime. Using the process of elimination, we can eliminate every option except for $\textbf{(A)}$ and $\textbf{(C)}$. Clearly, $606$ isn't prime, so the answer must be $\boxed{\textbf{(C) }2^{607}-1}$.

 Note: A organization of computer scientist and mathematicians named Great Internet Mersenne Prime Search (GIMPS) search for the worlds biggest prime numbers. More information can be found on \href{https://www.mersenne.org/primes/}{https://www.mersenne.org/primes/}~hashbrown2009

 



\textbf{Solution 3a (Elimination)}

We examine option E first. $2^{607}$ has a units digit of $8$ (Taking the units digit of the first few powers of two gives a pattern of $2, 4, 8, 6, 2, 4, 8, 6, 2, 4, 8, 6,\cdots$) and $3^{607}$ has a units digit of $7$ (Taking the units digit of the first few powers of three gives a pattern of $3, 9, 7, 1, 3, 9, 7, 1, 3, 9, 7, 1,\cdots$). Adding $7$ and $8$ together, we get $15$, which is a multiple of $5$, meaning that $2^{607}+3^{607}$ is divisible by 5.

 Next, we examine option D. We take the first few powers of $2$ added with $1$:
 \[2^1+1=3\]
 \[2^2+1=5\]
 \[2^3+1=9\]
 \[2^4+1=17\]
 \[2^5+1=33\]
 \[2^6+1=65\]
 \[2^7+1=129\]

 We see that the odd powers of $2$ added with 1 are multiples of three. If we continue this pattern, $2^{607}+1$ will be divisible by $3$. (The reason why this pattern works: When you multiply $2 \equiv2\pmod{3}$ by $2$, you obtain $4 \equiv1 \pmod{3}$. Multiplying by $2$ again, we get $1\cdot2\equiv2 \pmod{3}$. We see that in every cycle of two powers of $2$, it goes from $2 \pmod{3}$ to $1 \pmod{3}$ and back to $2 \pmod{3}$.)

 Next, we examine option B. We see that $2^{606}$ has a units of digits of $4$ (Taking the units digit of the first few powers of two gives a pattern of $2, 4, 8, 6, 2, 4, 8, 6, 2, 4, 8, 6,\cdots$). Adding $1$ to $4$, we get $5$. Since $2^{606}+1$ has a units digit of $5$, it is divisible by $5$.

 Lastly, we examine option A. Using the difference of cubes factorization $a^3-b^3=(a-b)(a^2+ab+b^2)$, we have $2^{606}-1^3=(2^{202}-1)(2^{404}+2^{202}+1)$. Since $2^{404}+2^{202}+1\equiv0\pmod{3}$ (Every term in the sequence is equivalent to $1\pmod{3}$), $2^{606}-1$ is divisible by $3$.

 Since we have eliminated every option except C, $\boxed{\text{(C)} \hspace{0.1 in}2^{607}-1}$ is not divisible by any prime less than $10$.

 



\textbf{Solution 3b (Elimination + Number Theory)}

We know that the prime numbers less than 10 are $2,3,5$ and $7$. We can start by testing if any of the answer choices are divisible by $2$. We see that they are all sums of one even number and one odd number, which is simply odd. So, we cannot exclude any answer choices so far. Now, let's check divisibility by $3$. We can use the fact that $2 \equiv -1 \pmod{3}$ to our advantage: 

 


 $(\text{A})\hspace{0.1in} 2^{606}-1$$\equiv (-1)^{606} -1$$\equiv 1-1$ $\equiv 0 \pmod{3}$
$(\text{B}) \hspace{0.1in} 2^{606}+1$$\equiv (-1)^{606} +1$$\equiv 1+1$$\equiv 2 \pmod{3}$
$(\text{C})\hspace{0.1in} 2^{607}-1$$\equiv (-1)^{607} -1$$\equiv -1 - 1$$\equiv -2 \pmod{3}$
$(\text{D})\hspace{0.1in} 2^{607} +1 \equiv (-1)^{607} +1 \equiv -1+1 \equiv 0 \pmod{3}$
$(\text{E}) \hspace{0.1in}  2^{607} + 3^{607} \equiv 2^{607} \equiv (-1)^{607} \equiv -1 \pmod{3}$

 
So, we eliminate choices A and D from divisibility by 3. Now, we move onto divisibility by 5. We can use cycling of powers to find useful remainders. Let's start with choice $B$. 

 


 $(\text{B}) \hspace{0.1in} 2^{606}+1$ 

 


 We see that the remainders of powers of $2$ when divided by $5$ cycle in a pattern: $2,4,3,1,2,4,3,1 \dots$. Since the pattern cycles every 4 terms, we use modulo 4 to simplify 606, getting that $606 \equiv 2\pmod{4}$. So, we get that $2^{606}$ is congruent modulo 5 to the second term of our pattern, so $2^{606} \equiv{4} \pmod{5}$. Thus, $2^{606}+1 \equiv 4+1 \equiv 5\equiv 0\pmod{5}$. We now eliminate choice B and compare choices C and E.  

 
$(\text{E}) \hspace{0.1in} 2^{607} + 3^{607}$ 

 


 
Looking at choice E, we see that we have to do similar cycling for powers of $3$. We get a pattern of $3,4,2,1,3,4,2,1 \dots$. Since this pattern also cycles every 4 terms, we use modulo 4 to simplify 607, getting that $607 \equiv 3\pmod{4}$. Both the power of 3 and the power of 2 have an exponent of 607, so we use the third term (since we just found that $607 \equiv 3\pmod{4}$) in each corresponding 4 term pattern to get that $2^{607} + 3^{607} \equiv 3+2 \equiv 5 \equiv 0\pmod{5}$. We eliminate choice E, and we are left with the correct answer: choice $\boxed{(\text{C}) \hspace{0.1in} 2^{607}-1}$




\textbf{Solution 3c (Elimination, Number Theory, and Function Analysis)}

Let’s start by calculating the unit’s digit of each answer choice.

 $2^{606}$ has a unit’s digit of 4, so choices A and B have unit’s digits of $3$ and $5$ respectively.

 $2^{607}$ has a unit’s digit of 8, so choices C and D have unit’s digits of $7$ and $9$ respectively.

 $3^{607}$ has a unit’s digit of 7. Combining with the statement above, choice E has a unit’s digit of $5$.

 So, none of the answer choices are divisible by $2$. But, we can eliminate choices B and E because they both have unit’s digit of $5$, and therefore, are divisible by $5$.

 We are left now with choices A, C, and D. Let’s now apply function analysis.

 Look at choices A and C again. They share the form $2^{n}-1$. Notice that if $n$ is even, it follows that $2^{n}-1$ is divisible by $3$. Since $606$ is even, choice A is divisible by $3$. However, our answer cannot be divisible by $3$, so eliminate A. This leaves us with choices C and D.

 Look at choice D again. It has the form $2^{n}+1$. Notice that if $n$ is odd, it follows that $2^{n}+1$ is divisible by 3. Since $607$ is odd, choice D is divisible by $3$. However, since the answer cannot be divisible by $3$, we eliminate D.

 At this point, you would choose choice C, but let’s now test divisibility by $7$. Notice that if $n$ is a multiple of 3, it follows that $2^{n}-1$ is divisible by $7$. However, since $607$ is not a multiple of $3$, choice C is not divisible by $7$.

 Since it is not divisible by any prime less than $10$, $\boxed{(\text{C}) \hspace{0.1in} 2^{607}-1}$ is the correct answer.

 



\textbf{Solution 4 (modulo 15 elimination)}

We proceed by taking the solutions modulo $15$, allowing us to eliminate both multiples of $3$ and multiples of $5$.

 First, we must identify cycles of powers of 2 modulo 15.$2^0 \equiv 1, \quad 2^1 \equiv 2, \quad 2^2 \equiv 4, \quad 2^3 \equiv 8, \quad 2^4 \equiv 1 \pmod{15}.$Thus, the pattern repeats every $4$ powers:$2^n \equiv \begin{cases} 1 &\text{if } n \equiv 0 \pmod{4},\\ 2 &\text{if } n \equiv 1 \pmod{4},\\ 4 &\text{if } n \equiv 2 \pmod{4},\\ 8 &\text{if } n \equiv 3 \pmod{4}. \end{cases}$

 Since $606 \equiv 2 \pmod{4}$ and $607 \equiv 3 \pmod{4}$,$2^{606} \equiv 4 \pmod{15}, \qquad 2^{607} \equiv 8 \pmod{15}.$

 Now let's look at cycles of powers of 3 mod 15.$3^0 \equiv 1, \quad 3^1 \equiv 3, \quad 3^2 \equiv 9, \quad 3^3 \equiv 12, \quad 3^4 \equiv 6, \quad 3^5 \equiv 3 \pmod{15}.$The cycle from here repeats every $4$ terms for $n \ge 1$:$3^n \equiv \begin{cases} 6 &\text{if } n \equiv 0 \pmod{4},\\ 3 &\text{if } n \equiv 1 \pmod{4},\\ 9 &\text{if } n \equiv 2 \pmod{4},\\ 12 &\text{if } n \equiv 3 \pmod{4}. \end{cases}$

 Since $607 \equiv 3 \pmod{4}$,$3^{607} \equiv 12 \pmod{15}$.

 We then evaluate each option modulo 15.

 \begin{aligned}\textbf{(A)} &:\; 2\textasciicircum{}{606}-1 \equiv 4-1 \equiv 3 \pmod{15},\\[6pt]\textbf{(B)} &:\; 2\textasciicircum{}{606}+1 \equiv 4+1 \equiv 5 \pmod{15},\\[6pt]\textbf{(C)} &:\; 2\textasciicircum{}{607}-1 \equiv 8-1 \equiv 7 \pmod{15},\\[6pt]\textbf{(D)} &:\; 2\textasciicircum{}{607}+1 \equiv 8+1 \equiv 9 \pmod{15},\\[6pt]\textbf{(E)} &:\; 2\textasciicircum{}{607}+3\textasciicircum{}{607} \equiv 8+12 \equiv 20 \equiv 5 \pmod{15}.\end{aligned}

 The residues divisible by $3$ or $5$ modulo $15$ are 0, 3, 5, 6, 9, 10, 12.Among our results, only $7$ (from option $\boxed{(\text{C})}$) is not among these, eliminating all the other answer choices.

 
$\boxed{(\text{C}) \hspace{0.1in} 2^{607}-1}$ is the number not divisible by any prime less than $10$. \\~Loquacious\_Autodidact

 



\textbf{Solution 5 (Contradiction; Less than 1 minute if you're good at number theory)}

This is like solution 1, but don't find the remainder, we assume it.

 The primes less than 10 are 2, 3, 5, and 7.

 Option choice $A$: Let $2^{606} \equiv 1 \pmod{3}$ as it clearly isn't divisible by 2. Then, because $2^2 \equiv 1 \pmod{3}$, it follows that $2^{606}$ satisfies the modulo. Option $A$ is not possible.

 Option choice $B$ ends in 5, and thus not possible.

 Option choice $C$: It is obviously not divisible by 2. We set $2^{607} \equiv 1 \pmod{3}$, which we see is not possible through Fermat's Little Theorem. We see $2^{607}$ doesn't end in 5. We set $2^{607} \equiv 1 \pmod{7}$, and through Fermat's Little Theorem $2^6 \equiv 1 \pmod{7}$. It follows that $607$ isn't a factor of 2, and thus does not hold the desired claim, a contradiction.

 At this point, for speed purposes, one would circle $C$ and move on. However, given the directory that this is a solution, I shall show the last two cases are not possible.

 Option choice $D$: It is obviously not divisible by 2. We set $2^{607} \equiv 2 \pmod{3}$. Through properties of modulo 3 we have that $(-1)^{607} \equiv -1 \pmod{3}$, which is true. Thus, $D$ is not plausible.

 Option choice $E$ is divisible by 5, and thus not possible. $\square$

 Additionally, for option $D$. One could notice that $2^{607} \equiv 2 \pmod{3}$ in the following way. We analyze the pattern. $2^1 \equiv 2 \pmod{3}$. $2^3 \equiv 2 \pmod{3}$. $2^5 \equiv 2 \pmod{3}$. $2^{2n+1} \pmod{3} \equiv 4^n \cdot 2^1 \pmod{3} \equiv 1^n \cdot 2 \pmod{3} \equiv 2 \pmod{3}$. 607 is odd, and thus $2^{607} \equiv 2 \pmod{3}$, as desired. $\square$




\textbf{Solution 6 (Extremely Illegal Cheese)}

All of the other solutions are extremely long, but this is how to cheese it using some methods that MIGHT be unreliable:First, turn 606 to 0 and 607 to 1. Next, Calculate each value, giving you $A=0 B=2 C=1 D=3 E=5$. Now since 2, 3, and 5 are prime, they can be divided by a prime(themselves). 0 doesn't really have multiples, so you can literally ignore it, and since 1 is only divisible by 1, which isnt prime, you get the correct answer, $\boxed{(\text{C})}$.

 


	\item (\textit{AMC 10 10B 2022}) Consider systems of three linear equations with unknowns $x$, $y$, and $z$,

 \begin{align*} a_1 x + b_1 y + c_1 z & = 0 \\ a_2 x + b_2 y + c_2 z & = 0 \\ a_3 x + b_3 y + c_3 z & = 0 \end{align*}

where each of the coefficients is either $0$ or $1$ and the system has a solution other than $x=y=z=0$.
For example, one such system is 
 \[\{ 1x + 1y + 0z = 0, 0x + 1y + 1z = 0, 0x + 0y + 0z = 0 \}\]
with a nonzero solution of $\{x,y,z\} = \{1, -1, 1\}$. How many such systems of equations are there?
(The equations in a system need not be distinct, and two systems containing the same equations in a
different order are considered different.)


 $\textbf{(A)}\ 302 \qquad\textbf{(B)}\ 338 \qquad\textbf{(C)}\ 340 \qquad\textbf{(D)}\ 343 \qquad\textbf{(E)}\ 344$

 \textbf{Answer: }B



\textbf{Solution 1 (Casework)}

Let $M_1=\begin{bmatrix}a_1 & b_1 & c_1\end{bmatrix}, M_2=\begin{bmatrix}a_2 & b_2 & c_2\end{bmatrix},$ and $M_3=\begin{bmatrix}a_3 & b_3 & c_3\end{bmatrix}.$ 

 We wish to count the ordered triples $(M_1,M_2,M_3)$ of row matrices. We perform casework:

 
            \begin{enumerate} 
        \item $M_1=M_2=M_3.$There are $2^3=8$ options for $M_1.$ Once $M_1$ is chosen, $M_2$ and $M_3$ are uniquely determined. 

 \textbf{In this case, we have $\boldsymbol{8}$ ordered triples $\boldsymbol{(M_1,M_2,M_3).}$}

  

 
        \item Exactly two of $M_1,M_2,$ and $M_3$ are equal.For $M_1=M_2\neq M_3,$ there are $2^3=8$ options for $M_1.$ Once $M_1$ is chosen, $M_2$ is uniquely determined, and $M_3$ has $2^3-1=7$ options. So, there are $8\cdot7=56$ ordered triples $(M_1,M_2,M_3).$

 Similarly, for each of $M_1=M_3\neq M_2$ and $M_2=M_3\neq M_1,$ there are $56$ ordered triples $(M_1,M_2,M_3).$

 \textbf{In this case, we have $\boldsymbol{56\cdot3=168}$ ordered triples $\boldsymbol{(M_1,M_2,M_3).}$}

  

 
        \item All of $M_1,M_2,$ and $M_3$ are different.There are two subcases: 

  

 <ol style="margin-left: 1.5em;" type="A"> 
        \item Exactly one of $M_1,M_2,$ and $M_3$ is $\begin{bmatrix}0 & 0 & 0\end{bmatrix}.$For $M_1=\begin{bmatrix}0 & 0 & 0\end{bmatrix},$ there are $2^3-1=7$ options for $M_2$ and $2^3-2=6$ options for $M_3.$ So, there are $7\cdot6=42$ ordered triples $(M_1,M_2,M_3).$

 Similarly, for each of $M_2=\begin{bmatrix}0 & 0 & 0\end{bmatrix}$ and $M_3=\begin{bmatrix}0 & 0 & 0\end{bmatrix},$ there are $42$ ordered triples $(M_1,M_2,M_3).$

 \textbf{In this subcase, we have $\boldsymbol{42\cdot3=126}$ ordered triples $\boldsymbol{(M_1,M_2,M_3).}$}

  

 
        \item The sum of two of $M_1,M_2,$ and $M_3$ is equal to the third matrix.For $M_1+M_2=M_3:$

            \begin{itemize} 
        \item If $M_1=\begin{bmatrix}1 & 0 & 0\end{bmatrix},$ then $M_2\in\{\begin{bmatrix}0 & 1 & 1\end{bmatrix},\begin{bmatrix}0 & 1 & 0\end{bmatrix},\begin{bmatrix}0 & 0 & 1\end{bmatrix}\}.$More generally, if $M_1$ consists of one $1$ and two $0$'s, then $M_2$ has $3$ options, and $M_3$ is uniquely determined. So, there are $\binom31\cdot3=9$ ordered triples $(M_1,M_2,M_3).$

        \item If $M_1=\begin{bmatrix}1 & 1 & 0\end{bmatrix},$ then $M_2=\begin{bmatrix}0 & 0 & 1\end{bmatrix}.$More generally, if $M_1$ consists of two $1$'s and one $0,$ then $M_2$ and $M_3$ are uniquely determined. So, there are $\binom32=3$ ordered triples $(M_1,M_2,M_3).$

            \end{itemize}
            There are $9+3=12$ ordered triples $(M_1,M_2,M_3).$ 

 Similarly, for each of $M_1+M_3=M_2$ and $M_2+M_3=M_1,$ there are $12$ ordered triples $(M_1,M_2,M_3).$

 \textbf{In this subcase, we have $\boldsymbol{12\cdot3=36}$ ordered triples $\boldsymbol{(M_1,M_2,M_3).}$}

  

 </ol>
            \end{enumerate}
            Together, the answer is $8+168+126+36=\boxed{\textbf{(B)}\ 338}.$




\textbf{Solution 2 (Intuitive, Straightforward)}

Begin by adding the three equations.
 \[(a_1 + a_2 + a_3)x + (b_1 + b_2 + b_3)y + (c_1 + c_2 + c_3)z = 0\].

 Each of the coefficients is 0 or 1, and $x, y, z$ cannot all equal 0. Notice that then, either $a_1 + a_2 + a_3 = 0, 1, 2, 3$, and the same for $b_1+ b_2 + b_3$, and $c_1 + c_2 + c_3$. Considering just $a_1 + a_2 + a_3$ for now, if equals to 0, then only 1 case exists. If equal to 1, then 3 cases exist. If equal to 2, then 3 cases still exist. If 3, then only 1 case exists. That is 8 total, so we have 24 total cases. Considering all 8 other viable systems (any permutation of $(x,y,z)$ besides $x =y = z = 0$), given the possible coefficient cases here, we have $24 \cdot 8 = 192$ possible cases.

 Now, subtract the three equations.
 \[(a_1-a_2-a_3)x - (b_1-b_2-b_3)y - (c_1-c_2-c_3)z= 0\]

 Now, notice that $a_1 - a_2 - a_3 = -2, -1, 0$, and the same for $b_1, b_2, b_3$ and $c_1, c_2, c_3$. For when $a_1 - a_2 -a_3 = -2$, that can only happen when we have only $(a_1, a_2, a_3) = (0, 1, 1)$, which is 1 way. If when $a_1 - a_2 - a_3 = -1$, we have this can happen in two ways, some permutations (except for when $a_1 = 1$) of either $(a_1, a_2, a_3) = (1,1,1), (0,0,1), (0,1,0)$, which is 3 ways total. For $a_1 - a_2 - a_3 = 0$, there are 3 ways, those being $(a_1,a_2,a_3) = (0,0,0), (1,0,1), (1,1,0)$. Finally, for $a_1 - a_2 - a_3 = 1$, we have the case $(a_1, a_2, a_3) = (1,0,0)$. Again 8 ways total, so times 3 is 24. When we consider this system, we only consider 6 plausible cases that do not degenerate into $x=y=z = 0$. To see where 6 comes from; due to subtraction, if $x+y = x+z = y+z = 0$, then subtracting will yield $x = y = z =0$ regardless, unlike the first case, where there could exist some $x$ or $y$ not necessarily equal to 0 but add to 0 that will still satisfy. Due to properties of subtraction, that is simply impossible for three distinct equations. So, all such equations must not contain any opposites of one another. Any subset from $\{-2, -1, -1, 0, 1, 1, 2 \}$ of three elements not consisting of both $-2$ and $2$ are viable roots. That is $\binom{7}{3} - 5 = 30$ cases. Only one fifth of these are viable as four fifths pertain values that will subtract to 0, a contradiction from the first claim (particularly, casework will show that, on $(-1,-1, z)$, $(1,1,z)$, $(0,-1,z)$, $(0,1,z)$, which is pertained by $4/5$ of the total will not work), so only 6 systems are valid. Therefore, we have $24 \cdot 6 = 144$ possible solutions.

 Now, the two forgotten solutions. When we add these two equations, we forget the case of redundancy (as we always count viable solutions, not infinite). When an equation is redundant, we must be able to deduce a viable ordered pair, which is the same as saying we must have systems in which their coefficients are of different parity. That is 2 cases. This however never applies for subtraction, as even if two equations have different parity, when we subtract the equations, that parity becomes constant, and will always cancel to 0, yielding only one possible solution, $(0,0,0)$, which is not allowed.

 So, exactly $192+144+2 = 336 + 2 = \boxed{338}$ solutions are viable

 



\textbf{Solution 3 (Solution 8 Alternate, Slightly Easier)}

We proceed with casework. The following will hold if 

 Case 1: At least one system has $0=0$, in which through the compliment, we have $512-343 = 169$ ways.

 Case 2: Exactly two equations are equal, in which we have $\binom{3}{2} \cdot 7 \cdot 6 = 126$ ways.

 Case 3: An equation is the result of the sum of the other two. The equation must have some coefficients between 0 and 1, so the possible cases for the coefficients of the third equation (excluding all 0 as that has already been counted) is $(0,0,1)$ or permutations, $(1,1,0)$ or permutations, and $(1,1,1)$. The first case has 0 ways. The second has 2 cases per each of the three cases, which is 6 ways total. The third has 6 cases in all, either the first equation is $(0,0,1)$ or some permutation (3 ways) which uniquely determines the second, or is (0,1,1) or some permutation (3 ways) which uniquely determines the second. So, there are $6 \cdot 6 = 36$ total cases.

 Case 4: All 3 are equal (except all 0, as already counted). That is 7 ways.

 So, exactly $169+126+36+7= \boxed{338}$ cases exist.

 



\textbf{Solution 4}

Let $n$V represent an equation in the system with $n$ variables. There are $2^9 = 512$ total possible systems (each coefficient can either be $0$ or $1$). We will subtract off the invalid ones (the equations with solution $x = y = z = 0$) using casework.

 We have $1$ case of selecting a 3-variable, 2-variable, and 1-variable equation, $6$ cases of selecting two $n$-variable equations and one different $m$-variable equation, and $3$ cases of selecting 3 of the same $m$-variable equations. However, we can't pick two 3-variable equations in our system because then it's impossible to get to the solution $x = y = z = 0.$ So there are $3$ invalid cases. Then we need to look for a total of $1 + 6 + 3 - 3 = 7$ different cases. (Doing this beforehand makes sure we don’t accidentally skip a case.)

 \textbf{Case 1: 3V-2V-2V:}WLOG say we have $x + y + z = 0$, $y + z = 0$, and $x + y = 0$. We are choosing two variables out of the 3 to be in the second equation, and we have 1 choice for the third. There are 6 ways to rearrange them. In total we get $6 \times \binom{3}{2} = 18$ systems for this case.

 \textbf{Case 2: 1V-1V-1V:}When we have $x = 0$, $y = 0$, and $z = 0$. There are 6 systems for this case.

 \textbf{Case 3: 2V-1V-1V:}We have 3 choices for the first equation: $x + y = 0$, $y + z = 0$, or $x + z = 0$. WLOG if $x + z = 0$ is the first equation, the second equation can either be $x = 0$ or $z = 0$, and the third equation must be $y = 0$. So we have 2 choices for the second equation and 1 choice for the third. $6 \times 3 \times 2$ gives us 36 systems for this case.

 \textbf{Case 4: 3V-2V-1V:} There is obviously one choice for the first equation ($x + y + z = 0$). For the second equation, we can choose any 2 of the 3 variables (in 3 ways). For the third, we must choose one of the two variables we included in our second equation (in 2 ways). In total, we get $6 \times 3 \times 2 = 36$ systems for this case.

 \textbf{Case 5: 2V-2V-2V:} There is only 1 combination of equations: $x + y = 0$, $y + z = 0$, and $x + z = 0$. These can be rearranged in 6 ways, giving us 6 systems.

 \textbf{Case 6: 3V-1V-1V:} Again, there's only one choice for the first equation ($x + y + z = 0$). For the second and third, we can choose any two of the three variables in 3 ways, giving us $6 \times 3 = 18$ systems.

 \textbf{Case 7: 2V-2V-1V:} We want to choose a variable to appear twice in both of the two variable equations. This can be done in 3 ways. The second equation will be the variable appearing twice added to the variable left out of the first (so only 1 choice), and 3 choices for the one variable equation. The system can be rearranged in 6 ways. This gives us $6 \times 3 \times 3 = 54$ systems.

 Our final answer is $512 - (18 + 6 + 36 + 36 + 6 + 18 + 54)  = \boxed{\textbf{(B)}\ 338}.$




\textbf{Solution 5 (Complementary Counting)}

We will use complementary counting and do casework on the equations.

 There are $8$ possible equations:

 Equation 1: $0 = 0$

 Equation 2: $x = 0$

 Equation 3: $y = 0$

 Equation 4: $z = 0$

 Equation 5: $x + y = 0$

 Equation 6: $x + z = 0$

 Equation 7: $y + z = 0$

 Equation 8: $x + y + z = 0$

 We will continue to refer to the equations by their number on this list.

 $8^3 = 512$ total systems. Note that no two equations by themselves can force $x = y = z = 0$. Therefore no system with Equation 1 or with repeated equations can force $x = y = z = 0$.

 \textbf{Case 1:} Equation 8 ($x + y + z = 0$) is present.

 Case 1a: Equation 8, and two equations from $\{5, 6, 7\}$.

 There are $\binom{3}{2} = 3$ ways to choose two equations from $\{5, 6, 7\}$ and $3! = 6$ ways to arrange each case. The number of options that force $x = y = z = 0$ is $3 \cdot 3! = 18$.

 Case 1b: Equation 8, one equation from $\{5, 6, 7\}$, and one equation from $\{2, 3, 4\}$.

 There are $\binom{3}{1} = 3$ ways to choose one equation from $\{5, 6, 7\}$. WLOG let us choose Equation 7. Given $x + y + z = 0$ and $y + z = 0$, we conclude that $x = 0$. The third equation can be either $y = 0$ or $z = 0$. There are $3!$ ways to arrange each case. The number of options that force $x = y = z = 0$ is $3 \cdot 2 \cdot 3! = 36$.

 Case 1c: Equation 8, and two equations from $\{2, 3, 4\}$.

 There are $\binom{3}{2} = 3$ ways to choose two equations from $\{2, 3, 4\}$ and $3! = 6$ ways to arrange each case. Each of these cases forces $x = y = z = 0$. $3 \cdot 3! = 18$ total options.

 \textbf{Case 2:} Equation 8 is $\textbf{not}$ present, at least one equation from $\{5, 6, 7\}$ is present.

 Case 2a: Equations $\{5, 6, 7\}$ are all present.

 There are $3!$ ways to arrange the three equations. $6$ options.

 Case 2b: Two equations from $\{5, 6, 7\}$ are present. One equation from $\{2, 3, 4\}$ is present.

 There are $\binom{3}{2}$ ways to choose two equations from $\{5, 6, 7\}$. WLOG let Equations 5 and 6 be in our system: $x + y = 0$ and $x + z = 0$. Any equation from $\{2, 3, 4\}$ will force $x = y = z = 0$. There are $3!$ ways to arrange the equations. The number of options that force $x = y = z = 0$ is $\binom{3}{2} \cdot \binom{3}{1} \cdot 3! = 54$.

 Case 2c: One equation from $\{5, 6, 7\}$ is present. Two equations from $\{2, 3, 4\}$ are present.

 There are $\binom{3}{1}$ ways to choose one equation from $\{5, 6, 7\}$. WLOG let Equation 5 ($x + y = 0$) be present. One of the two equations from $\{2, 3, 4\}$ must be Equation 4, $z = 0$, since it is the only equation that restricts $z$. The last equation can be either 2 or 3. There are $3!$ ways to arrange the equations. The number of options that force $x = y = z = 0$ is $\binom{3}{1} \cdot \binom{2}{1} \cdot 3! = 36$.

 \textbf{Case 3:} Only equations $\{2, 3, 4\}$ are present.

 There are $3!$ ways to arrange the three equations. $6$ options.

 We add up the cases: $18 + 36 + 18 + 6 + 54 + 36 + 6 = 174$ total systems force $x = y = z = 0$. Thus $512 - 174 = \boxed{\textbf{(B)}\ 338}$ do not.

 



\textbf{Solution 6 (Complementary Counting)}

The total number of possible systems is $2^9 = 512$, with $8$ possible sets of coefficients per equation. We will use complementary counting to find the number of systems which only have the solution $(0, 0, 0)$ and subtract that from the total. Similar to what is observed in Solution 3, if any equation is repeated or $0x + 0y + 0z = 0$, there will only be two or fewer equations for three variables, making one unique solution impossible. Therefore, we must choose $3$ different equations from $7$ possible ones, giving $7 \cdot 6 \cdot 5 = 210$ systems. However, there are two exceptions to consider, which will have more than one solution. The first is of the form $x + y + z = 0$, $x + y = 0$, $z = 0$; the second is of the form $x = 0$, $y = 0$, $x + y = 0$. In both cases, there are $3$ ways to choose the variables in the equations, and then $6$ ways to arrange them, giving $2 \cdot 3 \cdot 6 = 36$ exceptions. Subtracting this gives $210 - 36 = 174$ systems with only one solution, and the answer is then $512 - 174 = \boxed{\textbf{(B)}\ 338}$. 

 



\textbf{Solution 7 (Complementary Counting, Linear Dependence, Vector Analysis)}

Denote vector $\overrightarrow{i} = \left( i_1, i_2, i_3 \right)^T$ for $i \in \left\{ a, b, c \right\}$.Thus, we need to count how many vector tuples $\left( \overrightarrow{a} , \overrightarrow{b} , \overrightarrow{c} \right)$ are linearly dependent.

 We do complementary counting.

 First, the total number of vector tuples $\left( \overrightarrow{a} , \overrightarrow{b} , \overrightarrow{c} \right)$ is $\left( 2^3 \right)^3 = 512$.

 Second, we count how many many vector tuples $\left( \overrightarrow{a} , \overrightarrow{b} , \overrightarrow{c} \right)$ are linearly independent.

 To meet this condition, no vector can be a zero vector $\overrightarrow{0} = \left( 0, 0, 0 \right)^T$.

 Next, we do the casework analysis.

 Case $1^c$: Three vectors are all on axes.

 In this case, the number of $\left( \overrightarrow{a} , \overrightarrow{b} , \overrightarrow{c} \right)$ is $3!$.

 Case $2^c$: Two vectors are on axes and the third vector is not.

 We construct such an instance in the following steps.

 Step 1: We determine which two vectors lie on axes.

 The number of ways is $3$.

 Step 2: For two vectors selected in Step 1, we determine which two axes they lie on.

 The number of ways is $3 \cdot 2$.

 Step 3: For the third unselected vector, we determine its value.

 To make three vectors linear independent, the third vector cannot be on the plane formed by the first two vectors.So the number of ways is $3$.

 Following from the rule of product, the number of $\left( \overrightarrow{a} , \overrightarrow{b} , \overrightarrow{c} \right)$ in this case is $3 \cdot 3 \cdot 2 \cdot 3$.

 Case $3^c$: One vector is on an axis and the other two are not.

 We construct such an instance in the following steps.

 Step 1: We determine which vector lies on an axis.

 The number of ways is $3$.

 Step 2: For the selected vector, we determine which axis it lies on.

 The number of ways is $3$.

 Step 3: We determine the values of the two unselected vectors.

 First, to be linearly independent, these two vectors are distinct.Second, to be linearly independent, we cannot have one vector $(1,1,1)$ and another one that is a diagonal vector on the plane that is perpendicular to the first selected vector.

 Thus, the number or ways in this step is $4 \cdot 3-2 = 10$.

 Following from the rule of product, the number of $\left( \overrightarrow{a} , \overrightarrow{b} , \overrightarrow{c} \right)$ in this case is $3 \cdot 3 \cdot 10$.

 Case $(4.4)^c$: No vector is on any axis.

 In this case, any three distinct vectors are linearly independent.So the number of $\left( \overrightarrow{a} , \overrightarrow{b} , \overrightarrow{c} \right)$ in this case is $4 \cdot 3 \cdot 2$.

 Putting all cases together, the number of vector tuples $\left( \overrightarrow{a} , \overrightarrow{b} , \overrightarrow{c} \right)$ that are linearly independent is
 \[ 8^3 - 3!  - 3 \cdot 3 \cdot 2 \cdot 3 - 3 \cdot 3 \cdot 10 - 4 \cdot 3 \cdot 2 = \boxed{\textbf{(B)}\ 338}. \]~Steven Chen (Professor Chen Education Palace, www.professorchenedu.com)

 



\textbf{Solution 8 (Casework, but more intuitive)}

Case 1: At least one of the equations is $0=0$.By complementary counting, there are $512-343=\boxed{169}$ possible systems which satisfy this condition.

 Case 2: Exactly 2 of the equations are equal.By directly counting, we have $\tbinom{3}{2} \cdot 7 \cdot 6=\boxed{126}$, where there is $\tbinom{3}{2}$ places to put the two identical equations, and $7,6$ ways to choose the nonzero equations.

 Case 3: The sum of two different equations is equal to the last equation.The equation which serves as the total sum of the other two can either be $x+y,x+z,y+z,$ which each have $1$ way respectively to pick a set of equations, and $x+y+z$ has $3$ ways to pick a set of equations. Then, we have $6 \cdot 6=\boxed{36}$ for this case accounting for permutations.

 Case 4: All 3 of the equations are equal.There are $2^3-1=\boxed{7}$ systems which satisfy this case (they cannot all be $0$ because it's already counted in Case 1).

 Thus, there are $169+126+36+7=\boxed{\textbf{(B)}\ 338}$ total ways to pick such a system.

 


	\item (\textit{AMC 10 10B 2022}) Each square in a $5 \times 5$ grid is either filled or empty, and has up to eight adjacent neighboring squares, where neighboring squares share either a side or a corner. The grid is transformed by the following rules:


 
            \begin{itemize}
        \item Any filled square with two or three filled neighbors remains filled.
            \end{itemize}
            
            \begin{itemize}
        \item Any empty square with exactly three filled neighbors becomes a filled square.
            \end{itemize}
            
            \begin{itemize}
        \item All other squares remain empty or become empty.
            \end{itemize}
            A sample transformation is shown in the figure below.

                    \begin{figure}[H]
                        \centering
                        \includegraphics[width=0.5\textwidth]{Problem/AMC_10_10B_2022_19_Problem_Figure_1.png}
                    \end{figure}
                    
Suppose the $5 \times 5$ grid has a border of empty squares surrounding a $3 \times 3$ subgrid. How many initial configurations will lead to a transformed grid consisting of a single filled square in the center after a single transformation? (Rotations and reflections of the same configuration are considered different.)

                    \begin{figure}[H]
                        \centering
                        \includegraphics[width=0.5\textwidth]{Problem/AMC_10_10B_2022_19_Problem_Figure_2.png}
                    \end{figure}
                    
$\textbf{(A)}\ 14 \qquad\textbf{(B)}\ 18 \qquad\textbf{(C)}\ 22 \qquad\textbf{(D)}\ 26 \qquad\textbf{(E)}\ 30$

 \textbf{Answer: }C



\textbf{Solution}

There are two cases for the initial configuration:

 
            \begin{enumerate} 
        \item The center square is filled. Exactly two of the eight adjacent neighboring squares of the center are filled. Clearly, the only possibility is that the squares along one diagonal are filled, as shown below:

 
                    \begin{figure}[H]
                        \centering
                        \includegraphics[width=0.5\textwidth]{Solution/AMC_10_10B_2022_19_Solution_Figure_1.png}
                    \end{figure}
                    In this case, there are $2$ possible initial configurations. All rotations and reflections are considered.

  

 
        \item The center square is empty.Exactly three of the eight adjacent neighboring squares of the center are filled. The possibilities are shown below:

 
                    \begin{figure}[H]
                        \centering
                        \includegraphics[width=0.5\textwidth]{Solution/AMC_10_10B_2022_19_Solution_Figure_2.png}
                    \end{figure}
                    In this case, there are $4+4+4+8=20$ possible initial configurations. All rotations and reflections are considered.

 

 
            \end{enumerate}
            Together, the answer is $2+20=\boxed{\textbf{(C)}\ 22}.$




\textbf{Solution Alternate}

Key Idea:

 Only the $3$ x $3$ inner grid can initially contain filled squares, since the outer border is empty. We want exactly the center square to be filled after one step.

 Case 1: The center is initially filled.

 A filled square remains filled precisely when it has $2$ or $3$ filled neighbors. We examine the possible choices of filled neighbors of the center and require that no other square becomes filled.

 This gives $2$ valid configurations, counting rotations/reflections.

 Case 2: The center is initially empty.

 For the center to become filled, it must have exactly $3$ filled neighbors. There are several possible shapes for these three neighbors. Counting the configurations of each shape gives

 $2 + 4 + 4 + 4 + 8 = 22.$

 The total number of configurations is

 \begin{align*}2+(2+4+4+4+8)&=2+22\\&=\boxed{24}.\end{align*}


 Therefore the answer is $\boxed{\textbf{(C)}\ 22}.$

 \href{https://artofproblemsolving.com/wiki/index.php?title=User:Lunar\_paradox&action=edit&redlink=1}{Lunar paradox} (\href{https://artofproblemsolving.com/wiki/index.php?title=User\_talk:Lunar\_paradox&action=edit&redlink=1}{talk}) 18:06, 24 August 2026 (EDT)

 \textbf{Remark} This problem is based off of Conway's Game of Life, a cellular automaton created by British mathematician John Conway in 1970. The game follows the exact same pattern of rules as this problem, but is played on an infinite grid of cells.

 


	\item (\textit{AMC 10 10B 2022}) Let $ABCD$ be a rhombus with $\angle{ADC} = 46^{\circ}$. Let $E$ be the midpoint of $\overline{CD}$, and let $F$ be the point on $\overline{BE}$ such that $\overline{AF}$ is perpendicular to $\overline{BE}$. What is the degree measure of $\angle{BFC}$?


 $\textbf{(A)}\ 110 \qquad \textbf{(B)}\ 111 \qquad \textbf{(C)}\ 112 \qquad \textbf{(D)}\ 113 \qquad \textbf{(E)}\ 114$

 \textbf{Answer: }D



\textbf{Solution 1}

Extend segments $\overline{AD}$ and $\overline{BE}$ until they meet at point $G$.

 Because $\overline{AB} \parallel \overline{ED}$, we have $\angle ABG = \angle DEG$ and $\angle GDE = \angle GAB$, so $\triangle ABG \sim \triangle DEG$ by AA.

 Because $ABCD$ is a rhombus, $AB = CD = 2DE$, so $AG = 2GD$, meaning that $D$ is a midpoint of segment $\overline{AG}$.

 Now, $\overline{AF} \perp \overline{BE}$, so $\triangle GFA$ is right and median $FD = AD$.

 So now, because $ABCD$ is a rhombus, $FD = AD = CD$. This means that there exists a circle from $D$ with radius $AD$ that passes through $F$, $A$, and $C$.

 AG is a diameter of this circle because $\angle AFG=90^\circ$. This means that $\angle GFC = \angle GAC = \frac{1}{2} \angle GDC$, so $\angle GFC = \frac{1}{2}(180^\circ - 46^\circ)=67^\circ$, which means that $\angle BFC = \boxed{\textbf{(D)} \ 113}$




\textbf{Solution 3}

Let $\overline{AC}$ meet $\overline{BD}$ at $O$, then $AOFB$ is cyclic and $\angle FBO = \angle FAO$. Also, $AC \cdot BO = [ABCD] = 2 \cdot [ABE] = AF \cdot BE$, so $\frac{AF}{BO} = \frac{AC}{BE}$, thus $\triangle AFC \sim \triangle BOE$ by SAS, and $\angle OEB = \angle ACF$, then $\angle CFE = \angle EOC = \angle DAC = 67^\circ$, and $\angle BFC = \boxed{\textbf{(D)} \ 113}$

 A little bit faster: $AOFB$ is cyclic $\implies \angle OFE = \angle BAO$. 

 $AB \parallel CD \implies \angle BAO = \angle OCE$.

 Therefore $\angle OFE=\angle OCE \implies OECF$ is cyclic.

 Hence $\angle CFE=\angle COE=\angle CAD = 67^\circ$.

 



\textbf{Solution 4}

Observe that all answer choices are close to $112.5 = 90+\frac{45}{2}$. A quick solve shows that having $\angle D = 90^\circ$ yields $\angle BFC = 135^\circ = 90 + \frac{90}{2}$, meaning that $\angle BFC$ increases with $\angle D$. Substituting, $\angle BFC = 90 + \frac{46}{2} = \boxed{\textbf{(D)} \ 113}$.

 



\textbf{Solution 5 (Similarity and Circle Geometry)}

This solution refers to the \textbf{Diagram} section.

 We extend $AD$ and $BE$ to point $G$, as shown below:
                    \begin{figure}[H]
                        \centering
                        \includegraphics[width=0.5\textwidth]{Solution/AMC_10_10B_2022_20_Solution_Figure_1.png}
                    \end{figure}
                    We know that $AB=AD=2$ and $CE=DE=1$. 

 By AA Similarity, $\triangle ABG \sim \triangle DEG$ with a ratio of $2:1$. This implies that $2AD=AG$ and $AD \cong DG$, so $AG=2AD=2\cdot2=4$. That is, $D$ is the midpoint of $AG$.

 Note that as $\angle{AFG}$ has an angle of 90 deg and $AG=2DG$, we can redraw our previous diagram, but construct a circle with radius $AD$ or $2$ centered at $D$ and by extending $CD$ to point $H$, which is on the circle, as shown below:
                    \begin{figure}[H]
                        \centering
                        \includegraphics[width=0.5\textwidth]{Solution/AMC_10_10B_2022_20_Solution_Figure_2.png}
                    \end{figure}
                    Notice how $F$ and $C$ are on the circle and that $\angle CFE$ intercepts with $\overset{\Large\frown} {CG}$. 

 Let's call $\angle CFE = \theta$.

 Note that $\angle CDG$ also intercepts $\overset{\Large\frown} {CG}$, So $\angle CDG = 2\angle CFE$. 

 Let $\angle CDG = 2\theta$. Notice how $\angle CDG$ and $\angle ADC$ are supplementary to each other. We conclude that 
 \begin{align*} 2\theta &= 180-\angle ADC \\ 2\theta &= 180-46 \\ 2\theta &= 134 \\ \theta &= 67.  \end{align*}
Since $\angle BFC=180-\theta$, we have $\angle BFC=180-67=\boxed{\textbf{(D)} \ 113}$.

 Note: You can also find that CFE is half of CDG via circle theorems. We know CDG = 180 - 46 = 134, therefore making CFE 67 and BFC 113.

 



\textbf{Solution 6 (Simplification/Reduction)}

If angle $ADC$ was a right angle, it would be much easier. Thus, first pretend that $ADC$ is a right angle. $ABCD$ is now a square. WLOG, let each of the side lengths be 1. We can use the Pythagorean Theorem to find the length of line $AE$, which is $\sqrt{5}/2$. We want the measure of angle $BFC$, so to work closer to it, we should try finding the length of line $BF$. Angle $FAB$ and angle $ABF$ are complementary. Angle $ABF$ and angle $FBC$ are also complementary. Thus, $\sin FAB=\cos ABF=\sin FBC$. $\sin FAB=\sin FBC=(1/2)/(\sqrt{5}/2)=1/\sqrt{5}$. Since $\sin FAB=1\sqrt{5}$,and $AB=1$, $FB=\sin FAB$. It follows now that $FE=3*\sqrt{5}/10$. 

 Now, zoom in on triangle $BEC$. To use the Law of Cosines on triangle $FBC$, we need the length of $FC$. Use the Law of Cosines on triangle $EFC$. Cos $E=1/\sqrt{5}$. Thus, after using the Law of Cosines, $FC=\sqrt{2/5}$. 

 Since we now have SSS on $BEC$, we can get use the Law of Cosines. $\cos BFC=1/-\sqrt{2}$. $\arccos 1/-\sqrt{2}$ is 45, but if the cosine is negative that means that the angle is the supplement of the positive cosine value. $180-45=135$. Angle $BFC$ is $135^\circ$. 

 Realize that, around point F, there will always be 3 right angles, regardless of what angle $ADC$ is. There are only two angles that change when $ADC$ changes. Break up angle $BFC$ into angle $BFB'$, which is always 90 degrees, and angle $B'FC$, which we have discovered to to be half of $ADC$. Thus, when angle $ADC$ is 46 degrees, then $B'FC$ will be 23. $23+90=113$. Angle $BFC$ is $\boxed{\textbf{(D) }113}$ degrees.

 



\textbf{Solution 7}

Draw an auxiliary line from D to the midpoint of AB. Label it G. Then quadrilateral BGDE is a parallelogram. Hence AF and BD are perpendicular. Now, G being the midpoint of the hypontenuse of triangle ABF, it is the circumcenter of it. Thus, GF = GA and so DG is the perpendicular bisector of AF. Therefore, triangle AFD is also a isosceles triangle. Since AD = FD = CD, triangle CFD is also an isosceles triangle. Their one distinct angles' sum being 46 degrees, angle BFC = 113.

 Jeongha Cho

 


 



\textbf{Solution 8(lots and lots of tedious angle chasing)}

It is a well known fact that connecting the midpoints of the sides of a rhombus gives us a rectangle. We let the midpoint of side $DC$ as $E$, the midpoint of side $AD$ as $F$, the midpoint of side $AB$ as $G$, and the midpoint of side $BC$ as $H$. We can connect $E, F, G, H$ to get rectangle $EFGH$. Note that the obtuse angles of the rhombus are each $134$ degrees. We can perform a little bit of angle chasing following this diagram(I cannot draw diagrams in LaTex so the rest of this solution will be diagram-free). Let the intersection point between line segment $BE$ and line segment $GH$ as $I$. Let angle $FEB$ to be $\theta$. Performing more angle chasing leads us to finding that angle $GIB$ is angle $\theta$ and angle $IBG$ is $113 - \theta$. Let point $J$ be on $BE$ such that $AJ$ is perpendicular to $BE$. Then, by performing yet more angle chasing leads us to finding that angle $JEC$ is $90 - \theta + 23 = 113 - \theta$. We can predict that triangle $JEC$ is similar to triangle $BIG$. This is because since $GB$ and $EC$ both bisect the sides of the rhombus and one angle is common(the $113 - \theta$ angle). Therefore, we can safely say that all angles in these two triangles must be the same and thus angle $CJE = IGB = 67$ degrees. Thus, our desired angle, $CJB$ is simply just $180 - 67 = \boxed{113}$. 

 


 



\textbf{Solution 9 (educated guess)}

The answer choices must have something to do with 46. We can do some playing around, like trying $180 - 46 = 134$. Since the answers are all around the $110$s, we can start manipulating our $134$. We could divide by $2$, subtract $90$, etc. Dividing by $2$ gives $67$, and we can see $180 - 67 = 113$. However, we can also see that $180-((180-(90-46))/2) = 112$, but this seems much too complicated for a simple rhombus. Therefore, we can reasonably assume the answer is $\boxed{113}$. 

 


	\item (\textit{AMC 10 10B 2022}) Let $P(x)$ be a polynomial with rational coefficients such that when $P(x)$ is divided by the polynomial $x^2 + x + 1$, the remainder is $x + 2$, and when $P(x)$ is divided by the polynomial $x^2 + 1$, the remainder is $2x + 1$. There is a unique polynomial of least degree with these two properties. What is the sum of the squares of the coefficients of that polynomial?


 $\textbf{(A) } 10 \qquad \textbf{(B) } 13 \qquad \textbf{(C) } 19 \qquad \textbf{(D) } 20 \qquad \textbf{(E) } 23$

 \textbf{Answer: }E



\textbf{Solution 1 (Experimentation)}

Given that all the answer choices and coefficients are integers, we hope that $P(x)$ has positive integer coefficients.

 Throughout this solution, we will express all polynomials in base $x$. E.g. $x^2 + x + 1 = 111_{x}$.

 We are given:
 \[111a + 12 = 101b + 21 = P(x).\]We add $111$ and $101$ to each side and balance respectively:
 \[111(a - 1) + 123 = 101(b - 1) + 122 = P(x).\]We make the unit's digits equal:
 \[111(a - 1) + 123 = 101(b - 2) + 223 = P(x).\]We now notice that:
 \[111(a - 11) + 1233 = 101(b - 12) + 1233 = P(x).\]Therefore $a = 11_{x} = x + 1$, $b = 12_{x} = x + 2$, and $P(x) = 1233_{x} = x^3 + 2x^2 + 3x + 3$. $3$ is the minimal degree of $P(x)$ since there is no way to influence the $x$‘s digit in $101b + 21$ when $b$ is an integer. The desired sum is $1^2 + 2^2 +3^2+ 3^2 = \boxed{\textbf{(E)} \ 23}$

 P.S. The four computational steps can be deduced through quick experimentation.

 



\textbf{Solution 2}

Let $P(x) = Q(x)(x^2+x+1) + x + 2$, then $P(x) = Q(x)(x^2+1) + xQ(x) + x + 2$, therefore $xQ(x) + x + 2 \equiv 2x + 1 \pmod{x^2+1}$, or $xQ(x) \equiv x-1 \pmod{x^2+1}$. The rule says we can change any $x^2$ into $-1$. This is because $x^2 \equiv -1 \pmod{x^2+1}$ (since $x^2 - (-1) = x^2+1$, which is a multiple of $x^2+1$).To find the least possible degree $Q(x)$, it would be logical to first assume it's a constant. Let's say $Q(x)=c$, which makes $xQ(x) = cx$. Is $cx$ congruent to $x-1$? The difference is $cx - (x-1) = (c-1)x+1$. This is a polynomial of degree 1, so it's not a multiple of $x^2+1$ (a degree 2 polynomial). So, a constant won't work.Could $Q(x)$ be a simple degree 1 polynomial? Let's try $Q(x) = x+1$. Let's check if $Q(x) = x+1$ works. Calculate $xQ(x)$: $x(x+1) = x^2+x$. 

 Now, apply our rule: we can replace $x^2$ with $-1$. Therefore, $x^2+x$ becomes $-1+x$, which is $x-1$. This shows that $xQ(x)$ is congruent to $x-1$.Since we showed that no constant polynomial works, and we found a degree 1 polynomial that does work, the minimum degree must be 1. The simplest polynomial of degree 1 that satisfies the condition is $Q(x)=x+1$, and expanding gives $P(x) = x^3+2x^2+3x+3$. Summing the squares of coefficients gives $\boxed{\textbf{(E)} \ 23}$




\textbf{Solution 3}

Let $P(x) = (x^2+x+1)Q_1(x) + x + 2$, then $P(x) = (x^2+1)Q_1(x) + xQ_1(x) + x + 2$

 Also, $P(x) = (x^2+1)Q_2(x) + 2x + 1$

 We infer that $Q_1(x)$ and $Q_2(x)$ have same degree, we can assume $Q_1(x) = x + a$, and $Q_2(x) = x + b$, since $P(x)$ has least degree. If this cannot work, we will try quadratic, etc.

 Then we get:$(x^2+1)(Q_1(x) - Q_2(x)) + xQ_1(x) - x + 1 = 0$

 The constant term gives us:$(Q_1(x) - Q_2(x)) + 1 = 0$

 So $Q_1(x) - Q_2(x) = -1$

 Substituting this in gives: $-(x^2+1) + xQ_1(x) - x + 1 = 0$

 Solving this equation, we get$Q_1(x) = x + 1$ 

 Plugging this into our original equation we get $P(x) = x^3 + 2x^2 + 3x + 3$

 Verify this works with $P(x) = (x^2+1)Q_2(x) + 2x + 1$

 Therefore the answer is $1^2 + 2^2 + 3^2 + 3^2 = \boxed{\textbf{(E)} \ 23}$




\textbf{Solution 4 (Undetermined Coefficients)}

Notice that we cannot have the quotients equal to some constants, since the same constant will yield different constant terms for $P(x)$ (which is bad) and different constants will yield different first coefficients (also bad). Thus, we try setting the quotients equal to linear terms (for minimizing degree).

 Let $P(x)=(x^2+x+1)(ax+b)+(x+2)$ and $P(x)=(x^2+1)(ax+c)+(2x+1)$. The quotients have the same $x$ coefficient, since $P(x)$ must have the same $x^3$ coefficient in both cases. Expanding, we get 
 \[P(x)=ax^3+(a+b)x^2+(a+b+1)x+(b+2)\] and 
 \[P(x)=ax^3+cx^2+(a+2)x+(c+1).\]

 Equating coefficients, we get $b+2=c+1$, $a+b+1=a+2$, and $a+b=c$. From the second equation, we get $b=1$, then substituting into         the first, $c=2$. Finally, from $a+b=c$, we have $a=1$. Now, $P(x)=(x^2+x+1)(ax+b)+(x+2)=(x^2+x+1)(x+1)+(x+2)=x^3+2x^2+3x+3$ and our answer is 
 \[1^2+2^2+3^2+3^2=\boxed{\textbf{(E)} \ 23}.\]




\textbf{Solution 5 (Quick, but Not Quicker Than Solution 2)}

We construct the following equations in terms of $P(x)$ and the information given by the problem:
 \[\textbf{(1) } P(x)=(x^2+x+1)\cdot Q(x)+x+2\]
 \[\textbf{(2) } P(x)=(x^2+1)\cdot R(x)+2x+1\]Upon inspection, $Q(x)$ and $R(x)$ cannot be constant, so the smallest possible degree of $P(x)$ is $3,$ and both $Q(x)$ and $R(x)$ are linear. 

 Let $Q(x)=x-q$ and $R(x)=x-r.$ We know there will be values for $q$ and $r$ that make the below equation hold, so we can assume that $P(x)$ has a leading coefficient of $1$.

 Substituting these values in, and setting $\textbf{(1)}$ and $\textbf{(2)}$ equal to each  other,
 \[(x^2+x+1)(x-q)+x+2=(x^2+1)(x-r)+2x+1.\]We plug in $x=0$, yielding $r+1=q.$ Substituting this value into the above equation, 
 \[(x^2+x+1)(x-r-1)+x+2=(x^2+1)(x-r)+2x+1.\]Letting $x=1,$ we conclude that $r=-2,$ so $R(x)=x+2.$Therefore, 
 \[P(x)=(x^2+1)(x+2)+2x+1 = x^3+2x^3+3x+3.\]The requested sum is 
 \[1^2+2^2+3^2+3^2=\boxed{\textbf{(E) }23}\]

 -Benedict T (countmath1)

 



\textbf{Solution 6 (Similar to Solution 3)}

By remainder theorem, the polynomial can be written as follows.

 
 \[P(x) = (x^2+x+1)Q_{1}(x)+x+2 = (x^2+1)Q_{2}(x)+2x+1\]This is a timed exam, we can use the information given by answer choices. The answer choices tell us this is the polynomial with integer coefficients, and we need to find the polynomial with the least degree so we can assume both $Q_{1}(x)$ and $Q_{2}(x)$ are linear (the coefficient of x should be same).

 Then we can write $P(x)$ as a cubic polynomial.

 
 \[P(x) = (x^2+x+1)(ax+b)+x+2 = (x^2+1)(ax+c)+2x+1\]Substituting $x=0,1,-1$ to determine the value of $a$ and $b$.

 We have:
 \[P(0) = b+2 = c+1\]
 \[P(1) = 3a+3b+3 = 2a+2c+3\]
 \[P(-1) = -a+b+1 = -2a+2c-1\]

 We can solve the simultaneous equations: $a=1,b=1,c=2$

 Hence, $P(x)=(x^2+x+1)(x+1)+x+2=x^3+2x^2+3x+3$. The answer is $1^2+2^2+3^2+3^2=\boxed{\textbf{(E) }23}$




\textbf{Solution 7 (Bounding with Answer Choices)}

We are given: 

 $P(x) \equiv x + 2 \pmod{x^2+x+1}$
$P(x) \equiv 2x + 1 \pmod{x^2+1}$

 Since the problem asks us about the coefficients of the polynomial, we should evaluate $P(x)$ at $1$.

 Our congruences now become:

 $P(1) \equiv 1 + 2 \pmod{(1)^2+(1)+1}$ $\implies P(1) \equiv 3 \pmod{3} \implies P(1) \equiv 0 \pmod{3}$.

 $P(1) \equiv 2(1) + 1 \pmod{(1)^2+1}$ $\implies P(1) \equiv 3 \pmod{2} \implies P(1) \equiv 1 \pmod{2}$.

 We now know that the coefficients of the polynomial sum to an odd multiple of 3. 

 We now begin bounding the value of the sum of the coefficients squared. For our answer, the sum of the coefficients squared must be between $10$ and $23$. Trying a quadratic polynomial, if $P(1) = 3$, then the maximum of the coefficients squared is $9$ ($3^2+0^2+0^2$ is the maximum, any other permutation can be trivially shown to be lower). If $P(1) = 9$, the minimum of the coefficients squared is $27$ ($3^2+3^2+3^2$, once again, any other permutation can be trivially shown to be higher). Any higher value of $P(1)$ obviously will not work, so $P(x)$ cannot be a quadratic.

 Since a quadratic did not work, we now move up to a cubic. We start with $P(1) = 3$. Similarly to the quadratic, there is no way to get an answer within the range of the answer choices. 

 However, moving up to $P(1) = 9$, we get that the minimum of the coefficients squared is $81/4$ (which we get from all 4 coefficients being equal at $9/4$). Since this is the minimum, any other answers will be higher. We get that $\boxed{\textbf{(E) }23}$ is our answer, since it is the only answer choice greater than $81/4$.

 -SwordOfJustice

 



\textbf{Solution 8}

Since the divisors $x^2 + x + 1$ and $x^2 + 1$ are both powers of 2, and the remainders $x + 2$ and $2x + 1$ are both linear, we can assume that $P(x)$ has a minimum degree of 3.

 We can then set the coefficients of the polynomial as $P(x) = ax^3 + bx^2 + cx + d$.

 Then, using long division, we divide  $P(x) = ax^3 + bx^2 + cx + d$ by $x^2 + x + 1$ and obtain a remainder of $(c - b)x + (a + d - b)$, which is equal to $x + 2$

 We also divide $P(x) = ax^3 + bx^2 + cx + d$ by $x^2 + 1$ and obtain a remainder of $(c - a)x + (d - b)$, which is equal to $2x + 1$.

 From here we can set up a systems of equations:

 $c -b = 1$
$a + d - b = 2$
$c - a = 2$
$d - b = 1$

 Thus giving us the solutions $a = 1$, $b = 2$, $c = 3$, $d = 3$ and hence $P(x) = x^3 + 2x^2 + 3x + 3$

 Squaring the coefficients, then adding them, $1^2 + 2^2 + 3^2 + 3^2$ we get the final answer $\boxed{\textbf{(E)} \ 23}$




\textbf{Solution 9}

We assume $P(x)$ is a monic cubic polynomial. Now, we set $P(x)=x^{3}+ax^{2}+bx+c$, and set the equations:

 $x^{3}+ax^{2}+bx+c=Q_{1}(x)(x^{2}+x+1)+x+2$
$x^{3}+ax^{2}+bx+c=Q_{2}(x)(x^{2}+1)+2x+1$

 Plugging in $x=i$ for the second equation, we get $-i-a+bi+c=2i+1$, and matching coefficients gives $c-a=1$ and $b-1=2 \rightarrow b=3$. Now, we can guess values for $a$ and $c$ based on the answer choices. We have that the sum of the squares of the current known coefficients is $1^{2}+3^{2}=10$, and guessing $c=1, a=0$ clearly doesnt match up with the answer choices, so we guess $c=3, a=2$ and find that the answer is $\boxed{\textbf{(E)} \ 23}$




\textbf{Solution 10 (EZ CRT)}

If we let $x=10$, we just get a CRT problem. Note: this only works because all the coefficients are integers between 0 and 9. Our problem now is to find $n$ such that
 \[n\equiv 12\pmod {111}\]
 \[n\equiv 21\pmod {101}\]After using some careful Euclidean Algorithm (or just eyeing it), $n=1233$, so our answer is $1^2+2^2+3^2+3^2 = \boxed{\textbf{(E)} \ 23}$.

 



\textbf{Solution 11 (Polynomial CRT)}

$x^2 + x + 1$ and $x^2 + 1$ are relatively prime according to Bézout's theorem.

 $P(x) \equiv U(x) + V(x) \pmod{(x^2 + x + 1)(x^2 + 1)}$, where
 \[U(x) \equiv x + 2 \pmod{x^2 + x + 1}, \quad U(x) \equiv 0 \pmod{x^2 + 1}\]
 \[V(x) \equiv 2x + 1 \pmod{x^2 + 1}, \quad V(x) \equiv 0 \pmod{x^2 + x + 1}\]

 Notice that\begin{align}& x^2 + 1 \equiv -x \pmod{x^2 + x + 1}\\& (x^2 + 1)(x + 1) \equiv -x(x + 1) \equiv 1 \pmod{x^2 + x + 1}\\& U(x) \equiv -(x^2 + 1) + 2(x^2 + 1)(x + 1) \pmod{(x^2 + x + 1)(x^2 + 1)}\\& U(x) \equiv 2x^3 + x^2 + 2x + 1 \pmod{(x^2 + x + 1)(x^2 + 1)}\end{align}

 And\begin{align}& x^2 + x + 1 \equiv x \pmod{x^2 + 1}\\& (x^2 + x + 1)x \equiv x^2 \equiv -1 \pmod{x^2 + 1}\\& V(x) \equiv 2(x^2 + x + 1) - (x^2 + x + 1)x \pmod{(x^2 + x + 1)(x^2 + 1)}\\& V(x) \equiv -x^3 + x^2 + x + 2 \pmod{(x^2 + x + 1)(x^2 + 1)}\end{align}

 Therefore $P(x) \equiv U(x) + V(x) \equiv x^3 + 2x^2 + 3x + 3 \pmod{(x^2 + x + 1)(x^2 + 1)}$, the answer is $1^2 + 2^2 + 3^2 + 3^2 = \boxed{\textbf{(E)} \ 23}$



	\item (\textit{AMC 10 10B 2022}) Let $S$ be the set of circles in the coordinate plane that are tangent to each of the three circles with equations $x^{2}+y^{2}=4$, $x^{2}+y^{2}=64$, and $(x-5)^{2}+y^{2}=3$. What is the sum of the areas of all circles in $S$?


 $\textbf{(A)}~48\pi\qquad\textbf{(B)}~68\pi\qquad\textbf{(C)}~96\pi\qquad\textbf{(D)}~102\pi\qquad\textbf{(E)}~136\pi\qquad$

 \textbf{Answer: }E



\textbf{Solution 1}


                    \begin{figure}[H]
                        \centering
                        \includegraphics[width=0.5\textwidth]{Solution/AMC_10_10B_2022_22_Solution_Figure_1.png}
                    \end{figure}
                    The circles match up as follows: Case $1$ is brown, Case $2$ is blue, Case $3$ is green, and Case 4 is red.Let $x^2 + y^2 = 64$ be circle $O$, $x^2 + y^2 = 4$ be circle $P$, and $(x-5)^2 + y^2 = 3$ be circle $Q$.All the circles in S are internally tangent to circle $O$.There are four cases with two circles belonging to each:

 $*$ $P$ and $Q$ are internally tangent to $S$.

 $*$ $P$ and $Q$ are externally tangent to $S$.

 $*$ $P$ is externally and Circle $Q$ is internally tangent to $S$.

 $*$ $P$ is internally and Circle $Q$ is externally tangent to $S$.

 Consider Cases $1$ and $4$ together. Since circles $O$ and $P$ have the same center, the line connecting the center of $S$ and the center of $O$ will pass through the tangency point of both $S$ and $O$ and the tangency point of $S$ and $P$. This line will be the diameter of $S$ and have length $r_P + r_O = 10$. Therefore the radius of $S$ in these cases is $5$.

 Consider Cases $2$ and $3$ together. Similarly to Cases $1$ and $4$, the line connecting the center of $S$ to the center of $O$ will pass through the tangency points. This time, however, the diameter of $S$ will have length $r_P-r_O=6$. Therefore, the radius of $S$ in these cases is $3$.

 The set of circles $S$ consists of $8$ circles - $4$ of which have radius $5$ and $4$ of which have radius $3$. The total area of all circles in $S$ is $4(5^2\pi + 3^2\pi) = 136\pi \Rightarrow \boxed{\textbf{(E)}}$.

 -naman12

 



\textbf{Solution 2}

We denote by $C_1$ the circle that has the equation $x^2 + y^2 = 4$.We denote by $C_2$ the circle that has the equation $x^2 + y^2 = 64$.We denote by $C_3$ the circle that has the equation $(x-5)^2 + y^2 = 3$.

 We denote by $C_0$ a circle that is tangent to $C_1$, $C_2$ and $C_3$.We denote by $\left( u, v \right)$ the coordinates of circle $C_0$, and $r$ the radius of this circle.

 From the graphs of circles $C_1$, $C_2$, $C_3$, we observe that if $C_0$ is tangent to all of them, then $C_0$ must be internally tangent to $C_2$.We have
 \[ u^2 + v^2 = \left( 8 - r \right)^2 . \hspace{1cm} (1) \]

 We do the following casework analysis in terms of the whether $C_0$ is externally tangent to $C_1$ and $C_3$.

 Case 1: $C_0$ is externally tangent to $C_1$ and $C_3$.

 We have
 \[ u^2 + v^2 = \left( r + 2 \right)^2   \hspace{1cm} (2) \]and
 \[ (u-5)^2 + v^2 = \left( r + \sqrt{3} \right)^2 . \hspace{1cm} (3) \]

 Taking $(2) - (1)$, we get $r + 2 = 8 - r$. Thus, $r = 3$.We can further compute (omitted here) that there exist feasible $(u,v)$ with this given $r$.

 Case 2: $C_1$ is internally tangent to $C_0$ and $C_3$ is externally tangent to $C_0$.

 We have
 \[ u^2 + v^2 = \left( r - 2 \right)^2  \hspace{1cm} (2) \]and
 \[ (u-5)^2 + v^2 = \left( r + \sqrt{3} \right)^2 . \hspace{1cm} (3) \]

 Taking $(2) - (1)$, we get $r - 2 = 8 - r$. Thus, $r = 5$.We can further compute (omitted here) that there exist feasible $(u,v)$ with this given $r$.

 Case 3: $C_1$ is externally tangent to $C_0$ and $C_3$ is internally tangent to $C_0$.

 We have
 \[ u^2 + v^2 = \left( r + 2 \right)^2  \hspace{1cm} (2) \]and
 \[ (u-5)^2 + v^2 = \left( r - \sqrt{3} \right)^2 . \hspace{1cm} (3) \]

 Taking $(2) - (1)$, we get $r + 2 = 8 - r$. Thus, $r = 3$.We can further compute (omitted here) that there exist feasible $(u,v)$ with this given $r$.

 Case 4: $C_1$ is internally tangent to $C_0$ and $C_3$ is internally tangent to $C_0$.

 We have
 \[ u^2 + v^2 = \left( r - 2 \right)^2  \hspace{1cm} (2) \]and
 \[ (u-5)^2 + v^2 = \left( r - \sqrt{3} \right)^2 . \hspace{1cm} (3) \]

 Taking $(2) - (1)$, we get $r - 2 = 8 - r$. Thus, $r = 5$.We can further compute (omitted here) that there exist feasible $(u,v)$ with this given $r$.

 
Because the graph is symmetric with the $x$-axis, and for each case above, the solution of $v$ is not 0. Hence, in each case, there are two congruent circles whose centers are symmetric through the $x$-axis.

 Therefore, the sum of the areas of all the circles in $S$ is $2\left( 3^2 \pi +5^2 \pi +3^2 \pi +5^2 \pi \right) = \boxed{\textbf{(E) } 136 \pi}$.

 



\textbf{Solution 3}

The problem of counting the number of circles tangent to a given number of other circles, lines, or points is called Apollonius's Problem. Here, there are three circles; solutions to Apollonius's problem tell us that there are 8 possible tagent circles. When looking for all eight, we can notice that the diagram is symmetrical about the x-axis, so we only need to find four different circles and then reflect them to the other side. 

 After drawing in those circles, we can use any of the above methods to find the radii of the circle. The important step is to realize that the segment between tangent circles' centers must pass through the point of tangency. From here, we can prove that the center of the circle with equation $x^2 + y^2 = 4$ and the center of the circle tangent to all three must both lie on the radius of $x^2 + y^2 = 64$, allowing us to quickly find the radius of the tangent circle.

 For the tangent circles mostly above the x-axis, there are two of radius 5, and two of radius 3. The total area is $50 \pi + 18 \pi = 68 \pi$. As due to symmetry there are four other circles below the x-axis, the total area is $136 \pi$, so we choose $\boxed{ (E) }$. 

 


 


	\item (\textit{AMC 10 10B 2022}) Ant Amelia starts on the number line at $0$ and crawls in the following manner. For $n=1,2,3,$ Amelia chooses a time duration $t_n$ and an increment $x_n$ independently and uniformly at random from the interval $(0,1).$ During the $n$th step of the process, Amelia moves $x_n$ units in the positive direction, using up $t_n$ minutes. If the total elapsed time has exceeded $1$ minute during the $n$th step, she stops at the end of that step; otherwise, she continues with the next step, taking at most $3$ steps in all. What is the probability that Amelia’s position when she stops will be greater than $1$?


 $\textbf{(A) }\frac{1}{3} \qquad \textbf{(B) }\frac{1}{2} \qquad \textbf{(C) }\frac{2}{3} \qquad \textbf{(D) }\frac{3}{4} \qquad \textbf{(E) }\frac{5}{6}$

 \textbf{Answer: }C



\textbf{Solution 1 (Geometric Probability)}

Let $x$ and $y$ be random variables that are independently and uniformly distributed in the interval $(0,1).$ Note that 
 \[P(x+y\leq 1)=\frac{\frac12\cdot1^2}{1^2}=\frac12,\] as shown below:
                    \begin{figure}[H]
                        \centering
                        \includegraphics[width=0.5\textwidth]{Solution/AMC_10_10B_2022_23_Solution_Figure_1.png}
                    \end{figure}
                    Let $x,y,$ and $z$ be random variables that are independently and uniformly distributed in the interval $(0,1).$ Note that 
 \[P(x+y+z\leq 1)=\frac{\frac13\cdot\left(\frac12\cdot1^2\right)\cdot1}{1^3}=\frac16,\] as shown below:
                    \begin{figure}[H]
                        \centering
                        \includegraphics[width=0.5\textwidth]{Solution/AMC_10_10B_2022_23_Solution_Figure_2.png}
                    \end{figure}
                    We have two cases:

 
            \begin{enumerate} 
        \item Amelia takes exactly $2$ steps.We need $x_1+x_2>1$ and $t_1+t_2>1.$ So, the probability is 
 \[P(x_1+x_2>1)\cdot P(t_1+t_2>1)=\left(1-\frac12\right)\cdot\left(1-\frac12\right)=\frac14.\]

        \item Amelia takes exactly $3$ steps.We need $x_1+x_2+x_3>1$ and $t_1+t_2\leq1.$ So, the probability is 
 \[P(x_1+x_2+x_3>1)\cdot P(t_1+t_2\leq1)=\left(1-\frac16\right)\cdot\frac12=\frac{5}{12}.\]

            \end{enumerate}
            Together, the answer is $\frac14 + \frac{5}{12} = \boxed{\textbf{(C) }\frac{2}{3}}.$




\textbf{Solution 2 (Generalization and Induction)}

We can in fact find the probability that any number of randomly distributed numbers on the interval $[0, 1]$ sum to more than $1$ using geometric probability, as shown in the video below.

 If we graph the points that satisfy $x + y < 1$, $0 < x, y < 1$, we get the triangle with points $(0, 0)$, $(1, 0)$, and $(0, 1)$. If we graph the points that satisfy $x + y + z < 1$, $0 < x, y, z < 1$, we get the tetrahedron with points $(0, 0, 0)$, $(1, 0, 0)$, $(0, 1, 0)$, and $(0, 0, 1)$.

 Of course, the probability of either of these cases happening is simply the area/volume of the points we graphed divided by the total area of the graph, which is always $1$ (this would be much simpler than my calculus proof above).

 Thus, we can now solve for the probability that the sum is less than one for $n$ numbers using induction.

 $\textbf{Claim:}$ The probability that the sum is less than one is $\frac{1}{n!}$.

 $\textbf{Base Case:}$ For just $1$ number, the probability is $1$.

 $\textbf{Induction step:}$ Suppose that the probability for $n$ numbers is $\frac{1}{n!}$. We will prove that the probability for $n+1$ numbers is $\frac{1}{(n+1)!}$. To prove this, we consider that the area of an $n+1$-dimensional tetrahedron is simply the area/volume of the base times the height divided by $n+1$. 

 Of course, the area of the base is $\frac{1}{n!}$, and the height is $1$, and thus, we obtain $\frac{1}{n! \cdot (n+1)} = \frac{1}{(n+1)!}$ as our volume (this may be hard to visualize for higher dimensions). The induction step is complete.

 The probability of the sum being less than $1$ is $\frac{1}{n!}$, and the probability of the sum being more than $1$ is $\frac{n!-1}{n!}$. This trivializes the problem. The answer is 
 \[\frac{1}{2} \cdot \frac{2! - 1}{2!} + \frac{1}{2} \cdot \frac{3! - 1}{3!} = \boxed{\textbf{(C) }\frac{2}{3}}.\]




\textbf{Solution 3 (Observations)}

There are two cases: Amelia takes two steps or three steps. 

 The former case has a probability of $\frac{1}{2}$, as stated in Solution 1, and thus the latter also has a probability of $\frac{1}{2}$. 

 The probability that Amelia passes $1$ after two steps is also $\frac{1}{2}$, as it is symmetric to the probability above.

 Thus, if the probability that Amelia passes $1$ after three steps is $x$, our total probability is $\frac{1}{2} \cdot \frac{1}{2} + \frac{1}{2} \cdot x$. We know that $0 < x < 1$, and it is relatively obvious that $x > \frac{1}{2}$ (because the probability that $x > \frac{3}{2}$ is $\frac{1}{2}$). This means that our total probability is between $\frac{1}{2}$ and $\frac{3}{4}$, non-inclusive, so the only answer choice that fits is $\boxed{\textbf{(C) }\frac{2}{3}}$.

 \textbf{Remark (Calculus)} It is not immediately clear why three random numbers between $0$ and $1$ have a probability of $\frac{5}{6}$ of summing to more than $1$. Here is a proof:

 Let us start by finding the probability that two random numbers between $0$ and $1$ have a sum of more than $x$, where $0 \leq x \leq 1$. 

 Suppose that our two numbers are $y$ and $z$. Then, the probability that $y > x$ (which means that $y + z > x$) is $1 - x$, and the probability that $y < x$ is $x$.

 If $y < x$, the probability that $y + z > x$ is $1 - x + y$. This is because the probability that $y + z < x$ is equal to the probability that $z < x - y$, which is $x - y$, so our total probability is $1 - (x - y) = 1 - x + y$.

 Let us now find the average of the probability that $y + z > x$ when $y < x$. Since $y$ is a random number between $0$ and $x$, its average is $\frac{x}{2}$. Thus, our average is $1 - x + \frac{1}{2} = 1 - \frac{x}{2}$. 

 Hence, our total probability is equal to 
 \[1(1-x) + \left(1 - \frac{x}{2}\right)(x) = 1 - \frac{1}{2}x^2.\]Now, let us find the probability that three numbers uniformly distributed between $0$ and $1$ sum to more than $1$.

 Let our three numbers be $a$, $b$, and $c$. Then, the probability that $a + b + c > 1$ is equal to the probability that $b + c$ is greater than $1 - a$, which is equal to $1 - \frac{1}{2}(1 - a)^2$. 

 To find the total probability, we must average over all values of $a$. This average is simply equal to the area under the curve $1 - \frac{1}{2}(1-x)^2$ from $0$ to $1$, all divided by $1$. We can compute this value using integrals:
 \begin{align*} \frac{\int_0^1 \! 1 - \frac{1}{2}(1 - x)^2 \mathrm{d}x}{1} &= \int_0^1 \! 1 - \frac{1}{2}(1 - x)^2 \mathrm{d}x \\ &= 1 - \frac{1}{2}\int_0^1 \! (1 - x)^2 \mathrm{d}x \\ &= 1 - \frac{1}{2}\int_0^1 \! x^2 \mathrm{d}x \\ &= 1 - \frac{1}{2}\left(\frac{1}{3}\right) \\ &= \frac{5}{6}. \end{align*}
For those who don't know calculus, $\int_m^n \! f(x) \mathrm{d}x$ is the area under the curve $f(x)$ from $m$ to $n$.

 \textbf{Remark (Rigorous Calculus)} In the language of probability, we have three random variables $X$, $Y$, $Z$, each independent and uniformly distributed over the interval $(0, 1)$. We are interested in the probability that $X+Y+Z>1$, or $P(X+Y+Z>1)$.

 It follows from the definition of the joint probability density function $f(x, y, z)$ that (we assume that $0<x, y, z<1$):\begin{align}P(X+Y+Z>1)=\iiint_{x+y+z>1} f(x,y,z) \,dx\,dy\,dz.\end{align}Since $X, Y, Z$ are independent, we have $f(x, y, z)=f_X(x)*f_Y(y)*f_Z(z)$, where $f_X(x), f_Y(y), f_Z(z)$ represent the probability densities of $X, Y, Z$ respectively. Recall that $X, Y, Z$ are uniformly distributed over the interval $(0, 1)$. Hence, $f_X(x)=f_Y(y)=f_Z(z)=1$.

 It remains to evaluate the following triple integral\begin{align}\iiint_{x+y+z>1} \,dx\,dy\,dz=\frac{5}{6}.\end{align}Therefore, $P(X+Y+Z>1)=\frac{5}{6}$.

 respect to whoever would do this during the comp

 -guy\_who\_lieks\_frogs

 


	\item (\textit{AMC 10 10B 2022}) Consider functions $f$ that satisfy 
 \[|f(x)-f(y)|\leq \frac{1}{2}|x-y|\] for all real numbers $x$ and $y$. Of all such functions that also satisfy the equation $f(300) = f(900)$, what is the greatest possible value of

 \[f(f(800))-f(f(400))?\]
$\textbf{(A)}\ 25 \qquad\textbf{(B)}\ 50 \qquad\textbf{(C)}\ 100 \qquad\textbf{(D)}\ 150 \qquad\textbf{(E)}\ 200$

 \textbf{Answer: }B



\textbf{Solution 1 (Absolute Values and Inequalities)}

We have
 \begin{align*} |f(f(800))-f(f(400))| &\leq \frac12|f(800)-f(400)| &&(\bigstar) \\ &\leq \frac12\left|\frac12|800-400|\right| \\ &= 100, \end{align*}
from which we eliminate answer choices $\textbf{(D)}$ and $\textbf{(E)}.$

 Note that
 \begin{alignat*}{8} |f(800)-f(300)| &\leq \frac12|800-300| &&= 250, \\ |f(800)-f(900)| &\leq \frac12|800-900| &&= 50, \\ |f(400)-f(300)| &\leq \frac12|400-300| &&= 50, \\ |f(400)-f(900)| &\leq \frac12|400-900| &&= 250. \\ \end{alignat*}Let $a=f(300)=f(900).$ Together, it follows that 
 \begin{align*} |f(800)-a|&\leq 50, \\ |f(400)-a|&\leq 50. \\ \end{align*}
We rewrite $(\bigstar)$ as
 \begin{align*} |f(f(800))-f(f(400))| &\leq \frac12|f(800)-f(400)| \\ &= \frac12|(f(800)-a)-(f(400)-a)| \\ &\leq \frac12|50-(-50)| \\ &=\boxed{\textbf{(B)}\ 50}. \end{align*}
~MRENTHUSIASM

 



\textbf{Solution 2 (Lipschitz Condition)}

Denote $f(900)-f(600) = a$.Because $f(300) = f(900)$, $f(300) - f(600) = a$.

 Following from the Lipschitz condition given in this problem, $|a| \leq 150$ and
 \[ f(800) - f(600) \leq \min \left\{ a + 50 , 100 \right\} \]and
 \[ f(400) - f(600) \geq \max \left\{ a - 50 , -100 \right\} . \]Thus,
 \begin{align*} f(800) - f(400) & \leq \min \left\{ a + 50 , 100 \right\} - \max \left\{ a - 50 , -100 \right\}  \\ & = 100 + \min \left\{ a, 50 \right\} - \max \left\{ a , - 50 \right\} \\ & = 100 + \left\{ \begin{array}{ll} a + 50 & \mbox{ if } a \leq -50 \\ 0 & \mbox{ if } -50 < a < 50 \\ -a + 50 & \mbox{ if } a \geq 50 \end{array} \right. . \end{align*}
Thus, $f(800) - f(400)$ is maximized at $a = 0$, $f(800)-f(600) = 50$, $f(400)-f(600)=-50$, with the maximal value 100.

 By symmetry, following from an analogous argument, we can show that $f(800) - f(400)$ is minimized at $a = 0$, $f(800)-f(600) = -50$, $f(400)-f(600)=50$, with the minimal value $-100$.

 Following from the Lipschitz condition,
 \begin{align*} f(f(800)) - f(f(400)) & \leq \frac{1}{2} \left| f(800) - f(400) \right| \\ & \leq 50 . \end{align*}
We have already construct instances in which the second inequality above is augmented to an equality.

 Now, we construct an instance in which the first inequality above is augmented to an equality.

 Consider the following piecewise-linear function:
 \[ f(x) = \left\{ \begin{array}{ll} \frac{1}{2} \left( x - 300 \right) & \mbox{ if } x \leq 300 \\ -\frac{1}{2} \left( x - 300 \right) & \mbox{ if } 300 < x \leq 400 \\ \frac{1}{2} \left( x - 600 \right) & \mbox{ if } 400 < x \leq 800 \\ -\frac{1}{2} \left( x - 900 \right) & \mbox{ if } x > 800 \end{array} \right.. \]Therefore, the maximum value of $f(f(800)) - f(f(400))$ is$\boxed{\textbf{(B)}\ 50}$.

 



\textbf{Solution 3 (Slopes)}

Divide both sides by $|x - y|$ to get $\frac{|f(x) - f(y)|}{|x - y|} \leq \frac{1}{2}$. This means that when we take any two points on $f$, the absolute value of the slope between the two points is at most $\frac{1}{2}$.

 Let $f(300) = f(900) = c$, and since we want to find the maximum value of $|f(800) - f(400)|$, we can take the most extreme case and draw a line with slope $-\frac{1}{2}$ down from $f(300)$ to $f(400)$ and a line with slope $\frac{1}{2}$ up from $f(800)$ to $f(900)$. Then $f(400) = c - 50$ and $f(800) = c + 50$, so $|f(800) - f(400)| = |c + 50 - (c - 50)| = 100$, and this is attainable because the slope of the line connecting $f(400)$ and $f(800)$ still has absolute value less than $\frac{1}{2}$.

 Therefore, $|f(f(800)) - f(f(400))| \leq \frac{1}{2}|f(800) - f(400)| = \frac{1}{2}(100) = \boxed{\textbf{(B)}\ 50}$.

 



\textbf{Solution 4 (Translation)}

Consider $g(x) = f(x)-f(300)$. Then $g(x)$ satisfies all the conditions and $g(300) = g(900) = 0$. We would want $g(400)$ and $g(800)$ as distant from each other as possible. So assign $g(400) = -50$ and $g(800) = 50$, the possible lower and upper bounds respectively. It follows that one can obtain the upper bound for $|g(g(800)) - g(g(400))| = |g(50) - g(-50  )| \leq \frac{1}{2}(100) = \boxed{\textbf{(B)}\ 50}$ as the answer.

 


	\item (\textit{AMC 10 10B 2022}) Let $x_0,x_1,x_2,\dotsc$ be a sequence of numbers, where each $x_k$ is either $0$ or $1$. For each positive integer $n$, define 

 \[S_n = \sum_{k=0}^{n-1} x_k 2^k\]
Suppose $7S_n \equiv 1 \pmod{2^n}$ for all $n \geq 1$. What is the value of the sum  

 \[x_{2019} + 2x_{2020} + 4x_{2021} + 8x_{2022}?\]
$\textbf{(A) } 6 \qquad \textbf{(B) } 7 \qquad \textbf{(C) }12\qquad \textbf{(D) } 14\qquad \textbf{(E) }15$

 \textbf{Answer: }A



\textbf{Solution 1}

In binary numbers, we have 
 \[S_n = (x_{n-1} x_{n-2} x_{n-3} x_{n-4} \ldots x_{2} x_{1} x_{0})_2.\]It follows that 
 \[8S_n = (x_{n-1} x_{n-2} x_{n-3} x_{n-4} \ldots x_{2} x_{1} x_{0}000)_2.\]We obtain $7S_n$ by subtracting the equations:
 \[\begin{array}{clccrccccccr}   & (x_{n-1} & x_{n-2} & x_{n-3} & x_{n-4} & \ldots & x_2 & x_1 & x_0 & 0 & 0 & 0 \ )_2 \\ -\quad & & & & (x_{n-1} & \ldots & x_5 & x_4 & x_3 & x_2 & x_1 & x_0)_2 \\ \hline   & & & & & & & & & & &  \\ [-2.5ex]   & ( \ \ ?& ? & ? & 0 \ \ \ & \ldots & 0 & 0 & 0 & 0 & 0 & 1 \ )_2 \\ \end{array}\]We work from right to left:
 \begin{alignat*}{6} x\_0=x\_1=x\_2=1  \quad &\implies \quad &x\_3 &= 0& \\  \quad &\implies \quad &x\_4 &= 1& \\  \quad &\implies \quad &x\_5 &= 1& \\ \quad &\implies \quad &x\_6 &= 0& \\  \quad &\implies \quad &x\_7 &= 1& \\  \quad &\implies \quad &x\_8 &= 1& \\ \quad &\quad \vdots & & & \end{alignat*}For all $n\geq3,$ we conclude that

 
            \begin{itemize}
        \item $x_n=0$ if and only if $n\equiv 0\pmod{3}.$
            \end{itemize}
            
            \begin{itemize}
        \item $x_n=1$ if and only if $n\not\equiv 0\pmod{3}.$
            \end{itemize}
            Finally, we get $(x_{2019},x_{2020},x_{2021},x_{2022})=(0,1,1,0),$ from which 
 \[x_{2019} + 2x_{2020} + 4x_{2021} + 8x_{2022} = \boxed{\textbf{(A) } 6}.\]~Steven Chen (Professor Chen Education Palace, www.professorchenedu.com)

 



\textbf{Solution 2}

First, notice that 
 \[x_{2019} + 2x_{2020} + 4x_{2021} + 8x_{2022} = \frac{S_{2023} - S_{2019}}{2^{2019}}.\]Then since $S_n$ is the modular inverse of $7$ in $\mathbb{Z}_{2^n}$, we can perform the Euclidean Algorithm to find it for $n = 2019,2023$.

 Starting with $2019$, 
 \begin{align*} 7S_{2019} &\equiv 1 \pmod{2^{2019}} \\ 7S_{2019} &= 2^{2019}k + 1. \end{align*}
Now, take both sides $\operatorname{mod} \ 7$:
 \[0 \equiv 2^{2019}k + 1 \pmod{7}.\]Using Fermat's Little Theorem, 
 \[2^{2019} = (2^{336})^6 \cdot 2^3 \equiv 2^3 \equiv 1 \pmod{7}.\]Thus,  
 \[0 \equiv k + 1 \pmod{7} \implies k \equiv 6 \pmod{7} \implies k = 7j + 6.\]Therefore, 
 \[7S_{2019} = 2^{2019} (7j + 6) + 1 \implies S_{2019} = \frac{2^{2019} (7j + 6) + 1}{7}.\]

 We may repeat this same calculation with $S_{2023}$ to yield 
 \[S_{2023} = \frac{2^{2023} (7h + 3) + 1}{7}.\]Now, we notice that $S_n$ is basically an integer expressed in binary form with $n$ bits.This gives rise to a simple inequality, 
 \[0 \leqslant S_n \leqslant 2^n.\]Since the maximum possible number that can be generated with $n$ bits is 
 \[\underbrace{{11111\dotsc1}_2}_{n} = \sum_{k=0}^{n-1} 2^k = 2^n - 1 \leqslant 2^n.\]Looking at our calculations for $S_{2019}$ and $S_{2023}$, we see that the only valid integers that satisfy that constraint are $j = h = 0$.
 \[\frac{S_{2023} - S_{2019}}{2^{2019}} = \frac{\frac{2^{2023} \cdot 3 + 1}{7} - \frac{2^{2019} \cdot 6 + 1}{7}}{2^{2019}} = \frac{2^4 \cdot 3 - 6}{7} = \boxed{\textbf{(A) } 6}.\]~zoomanTV

 



\textbf{Solution 3}

As in Solution 2, we note that 
 \[x_{2019}+2x_{2020}+4x_{2021}+8x_{2022}=\frac{S_{2023}-S_{2019}}{2^{2019}}.\]We also know that $7S_{2023} \equiv 1 \pmod{2^{2023}}$ and $7S_{2019} \equiv 1 \pmod{2^{2019}}$, this implies:
 \[\textbf{(1) } 7S_{2023}=2^{2023}\cdot{x} + 1,\]
 \[\textbf{(2) } 7S_{2019}=2^{2019}\cdot{y} + 1.\]Dividing by $7$, we can isolate the previous sums:
 \[\textbf{(3) } S_{2023}=\frac{2^{2023}\cdot{x} + 1}{7},\]
 \[\textbf{(4) } S_{2019}=\frac{2^{2019}\cdot{y} + 1}{7}.\]The maximum value of $S_n$ occurs when every $x_i$ is equal to $1$. Even when this happens, the value of $S_n$ is less than $2^n$. Therefore, we can construct the following inequalities:
 \[\textbf{(3) } S_{2023}=\frac{2^{2023}\cdot{x} + 1}{7} < 2^{2023},\]
 \[\textbf{(4) } S_{2019}=\frac{2^{2019}\cdot{y} + 1}{7} < 2^{2019}.\]From these two equations, we can deduce that both $x$ and $y$ are less than $7$. 

 Reducing $\textbf{1}$ and $\textbf{2}$ $\pmod{7},$ we see that 
 \[2^{2023}\cdot{x}\equiv 6\pmod{7},\]and 
 \[2^{2019}\cdot{y}\equiv 6\pmod{7}.\]

 The powers of $2$ repeat every $3, \pmod{7}.$ 

 Therefore, $2^{2023}\equiv 2 \pmod 7$ and $2^{2019} \equiv 1 \pmod {7}.$ Substituing this back into the above equations, 
 \[2x\equiv{6}\pmod{7}\]and
 \[y\equiv{6}\pmod{7}.\]

 Since $x$ and $y$ are integers less than $7$, the only values of $x$ and $y$ are $3$ and $6$ respectively. 

 The requested sum is 
 \begin{align*} \frac{S_{2023}-S_{2019}}{2^{2019}} &= \frac{\frac{2^{2023}\cdot{x} + 1}{7} - \frac{2^{2019}\cdot{y} + 1}{7}}{2^{2019}} \\ &= \frac{1}{2^{2019}}\left(\frac{2^{2023}\cdot{3} + 1}{7} -\left(\frac{2^{2019}\cdot{6} + 1}{7}  \right)\right) \\ &= \frac{3\cdot{2^4}-6}{7} \\ &= \boxed{\textbf{(A) } 6}. \end{align*}
-Benedict T (countmath1)

 



\textbf{Solution 4}

Note that, as in Solution 2, we have 
 \[x_{2019} + 2x_{2020} + 4x_{2021} + 8x_{2022} = \frac{S_{2023} - S_{2019}}{2^{2019}}.\] This is because 
 \[S_{2023} = x_{0}2^{0} + x_{1}2^{1} + \cdots + x_{2019}2^{2019} + \cdots + x_{2022}2^{2022}\] and 
 \[S_{2019} = x_{0}2^{0} + x_{1}2^{1} + \cdots + x_{2018}2^{2018}.\] Note that 
 \[S_{2023} - S_{2019} = x_{2019}2^{2019} + \cdots + x_{2022}2^{2022} = 2^{2019}(x_{2019} + 2x_{2020} + 4x_{2021} + 8x_{2022}).\] Therefore, 
 \[x_{2019} + 2x_{2020} + 4x_{2021} + 8x_{2022} = \frac{S_{2023} - S_{2019}}{2^{2019}}.\] Multiplying both sides by 7 gives us 
 \[7(x_{2019} + 2x_{2020} + 4x_{2021} + 8x_{2022}) = \frac{7S_{2023} - 7S_{2019}}{2^{2019}}.\] We can write 
 \[7S_{2023} = 1\pmod{2^{2023}} = 1 + 2^{2023}a = 1 + 2^{2019}*16a\] and 
 \[7S_{2019} = 1\pmod{2^{2019}} = 1 + 2^{2019}b\] for some a and b. Substituting, we get 
 \[7(x_{2019} + 2x_{2020} + 4x_{2021} + 8x_{2022}) = \frac{(1 + 2^{2019} * 16a) - (1 + 2^{2019}b)}{2^{2019}} = 16a - b.\] Therefore, our answer can be written as 
 \[\frac{16a - b}{7}.\] Another thing to notice is that a and b are integers between 0 and 6. This is because 
 \[7(1 + 2 + 4 + 8 + \cdots + 2^{2022}) \geqslant 7S_{2023} = 1 + 2^{2023}a\] which is 
 \[7(2^{2023}) - 7 \geqslant 1 + 2^{2023}a\] 
 \[(7-a) \geqslant \frac{8}{2^{2023}}\] which only holds when a is less than 7 because the right is very small positive number, so the left must be positive, too. Clearly, a is also non-negative, because otherwise, 
 \[7S_{2023} = 1 + 2^{2023}a < 0\] which would mean 
 \[S_{2023} < 0\] which cannot happen, so a is greater than 0.  A similar explanation for b shows that b is an integer between 0 and 6 inclusive. 

 Going back to the solution, if our answer to the problem is n, then 
 \[16a - b = 7n\] and 
 \[16a = 7n + b,\] so we can try the five option choices and see which one, when multiplied by 7 and added to some whole number between 0 and 6 results in a multiple of 16. Trying all the option choices, we see that you need to add 7n to something more than 6 to equal a multiple of 16 other than for option A. Therefore, the answer is $\boxed{\textbf{(A) } 6}.$



\end{enumerate}